\documentclass[12pt,a4paper]{report}

% ---------------------------------------------------------------------
% Codificação, idioma, fontes
% ---------------------------------------------------------------------
\usepackage[utf8]{inputenc}
\usepackage[T1]{fontenc}
\usepackage[brazil]{babel}
\usepackage{lmodern}
\usepackage{microtype}

% ---------------------------------------------------------------------
% Matemática
% ---------------------------------------------------------------------
\usepackage{amsmath}
\usepackage{amssymb}
\usepackage{amsthm}
\usepackage{bm}
\usepackage{siunitx}
\sisetup{output-decimal-marker = {,}, per-mode = symbol}

% ---------------------------------------------------------------------
% Tabelas, figuras, layout
% ---------------------------------------------------------------------
\usepackage{graphicx}
\usepackage{booktabs}
\usepackage{longtable}
\usepackage{array}
\usepackage{float}
\usepackage[margin=2.5cm]{geometry}
\usepackage{setspace}
\onehalfspacing
\usepackage{enumitem}
\usepackage{xcolor}

% ---------------------------------------------------------------------
% Diagramas
% ---------------------------------------------------------------------
\usepackage{tikz}
\usetikzlibrary{shapes.geometric, arrows.meta, positioning, fit, backgrounds, calc, decorations.pathreplacing}

% ---------------------------------------------------------------------
% Código-fonte
% ---------------------------------------------------------------------
\usepackage{listings}
\lstdefinelanguage{Fortran90}{
  morekeywords={program,module,use,implicit,none,subroutine,function,end,do,end do,if,then,else,end if,call,integer,real,character,allocatable,intent,in,out,inout,dimension,type,contains,return,stop,parameter,logical,allocate,deallocate,where,elsewhere},
  sensitive=false,
  morecomment=[l]{!},
  morestring=[b]",
}
\lstset{
  language=Fortran90,
  basicstyle=\ttfamily\small,
  keywordstyle=\color{blue!70!black}\bfseries,
  commentstyle=\color{green!40!black}\itshape,
  stringstyle=\color{orange!70!black},
  numbers=left,
  numberstyle=\tiny\color{gray},
  frame=single,
  breaklines=true,
  breakatwhitespace=false,
  showstringspaces=false,
  tabsize=2,
  columns=flexible,
}

% ---------------------------------------------------------------------
% Referências cruzadas e hyperlinks
% ---------------------------------------------------------------------
\usepackage[
  backend=biber,
  style=numeric,
  sorting=nyt,
  maxbibnames=6,
]{biblatex}
\addbibresource{referencias.bib}

\usepackage[hidelinks,
  colorlinks=true,
  linkcolor=black,
  citecolor=blue!50!black,
  urlcolor=blue!50!black,
]{hyperref}
\usepackage[capitalize]{cleveref}

% ---------------------------------------------------------------------
% Cabeçalhos
% ---------------------------------------------------------------------
\usepackage{fancyhdr}
\pagestyle{fancy}
\fancyhf{}
\fancyhead[LE,RO]{\thepage}
\fancyhead[RE]{\nouppercase{\leftmark}}
\fancyhead[LO]{\nouppercase{\rightmark}}
\renewcommand{\headrulewidth}{0.4pt}

% ---------------------------------------------------------------------
% Ambientes auxiliares
% ---------------------------------------------------------------------
\newtheorem{definicao}{Definição}[chapter]
\newtheorem{achado}{Achado técnico}[chapter]
\newenvironment{caixaachado}
  {\begin{achado}\begin{quote}\itshape}
  {\end{quote}\end{achado}}

\newcommand{\selfgrib}{\textit{selfgrib}}
\newcommand{\mpasa}{MPAS-A}

% ---------------------------------------------------------------------
% Metadados
% ---------------------------------------------------------------------
\title{
  \Large\textbf{Interpolação Nativa em Malha de Voronoi para o Modelo MPAS-A:\\[0.3em]
  Geração de Condições Inicial e de Fronteira Lateral\\Sem Reamostragem em Grade Regular}
  \\[1em]
  \large Relatório técnico-científico do desenvolvimento \emph{selfgrib} / rota \emph{voronoi-to-voronoi}
}
\author{Diego Pereira \\ \small Instituto Nacional de Pesquisas Espaciais (INPE) \\ \small \texttt{diego.pereira@inpe.br}}
\date{Setembro de 2026}

\begin{document}

\begin{titlepage}
\maketitle
\thispagestyle{empty}
\vfill
\begin{abstract}
\noindent
Este relatório documenta o desenvolvimento de uma rota alternativa de
geração de condições inicial (\texttt{init.nc}) e de fronteira lateral
(\texttt{lbc.*.nc}) para o núcleo de previsão \mpasa{} (\textit{Model for
Prediction Across Scales -- Atmosphere}) que interpola diretamente entre
duas malhas de Voronoi esféricas -- a malha nativa de uma rodada global
já concluída e a malha regional de destino -- sem nenhuma etapa
intermediária de reamostragem em grade latitude-longitude regular.
O trabalho parte do \selfgrib{}, um \textit{ungrib} alternativo já
existente que usa a saída nativa de uma rodada global do próprio
\mpasa{} como fonte de dado meteorológico, eliminando a dependência de
GRIB externo (GFS, BAM, Eta). A rota descrita aqui vai um passo além:
em vez de reescrever o formato binário intermediário do WPS e deixar o
\texttt{init\_atmosphere\_model} original reinterpolar esse dado de volta
para a malha nativa (usando bilinear/16-pontos sobre uma grade
estruturada -- o único modo de interpolação que esse programa aceita),
o pipeline aqui apresentado reimplementa, em Fortran, extraindo
literalmente as fórmulas do código-fonte de referência do MPAS-Model
(versão 3.0.2), cada etapa necessária -- interpolação horizontal
baricêntrica no triângulo dual de Delaunay, geração da grade vertical
terrain-following híbrida, interpolação vertical, balanço hidrostático,
inicialização de campos de superfície/solo -- escrevendo os arquivos
\texttt{init.nc}/\texttt{lbc.*.nc} diretamente, prontos para alimentar o
núcleo dinâmico \texttt{mpas\_atmosphere} sem qualquer grade intermediária.
A implementação foi validada campo a campo contra dados reais de
produção (malha regional \texttt{SouthAmerica}, recorte de
$x1.163842$, $\sim\SI{60}{\kilo\metre}$ de resolução) e, por fim, contra
o próprio consumidor final: uma previsão de 24 horas rodada com o
\texttt{mpas\_atmosphere} real a partir dos arquivos gerados por essa
rota, sem erros e com campos fisicamente sãos, cuja divergência frente à
rota de referência (via WPS) cresce de forma suave e limitada ao longo do
tempo de integração -- o padrão esperado de duas trajetórias vizinhas de
um sistema caótico, não de um artefato numérico.

\vspace{1em}
\noindent\textbf{Palavras-chave:} MPAS-A; malha de Voronoi; triangulação
de Delaunay; interpolação baricêntrica; coordenada vertical
\textit{terrain-following}; condição inicial; fronteira lateral;
previsão numérica do tempo.
\end{abstract}
\vfill
\end{titlepage}

\tableofcontents
\listoffigures
\listoftables
\clearpage

\chapter{Introdução}
\label{cap:introducao}

\section{Motivação}

O \mpasa{} (\textit{Model for Prediction Across Scales -- Atmosphere})
\cite{skamarock2012voronoi} é um modelo global/regional de previsão
numérica do tempo cujo núcleo dinâmico opera sobre uma malha horizontal
não estruturada: uma tesselação de Voronoi centroidal esférica (SCVT),
cujos polígonos (predominantemente hexágonos, com um pequeno número de
pentágonos e heptágonos nos pontos de maior curvatura da geração da
malha) definem os volumes de controle das variáveis de massa, enquanto
as arestas entre células vizinhas carregam a componente normal do vento.
Essa escolha de malha -- em contraste com as grades latitude-longitude
regulares tradicionalmente usadas em modelos globais mais antigos --
evita a convergência de meridianos nos polos e permite refinamento local
suave sem as descontinuidades de grades aninhadas.

O projeto no qual este trabalho se insere, \selfgrib{}, nasceu de uma
pergunta simples: é possível gerar a condição inicial e a fronteira
lateral de uma rodada regional do \mpasa{} usando, como fonte de dado
meteorológico, a saída nativa de uma rodada global do \emph{próprio}
\mpasa{}, sem depender de nenhum GRIB externo (GFS, BAM, Eta) e sem
aplicar nenhum \emph{patch} no código-fonte do MPAS-Model? A resposta,
documentada em detalhe no relatório principal do projeto, foi sim: o
\selfgrib{} lê o \texttt{history.nc} de uma rodada global concluída,
extrai e interpola os campos necessários, e escreve o formato binário
intermediário do WPS -- o mesmo formato que o \texttt{ungrib} geraria a
partir de um GRIB real -- permitindo que o \texttt{init\_atmosphere\_model}
original consuma esse dado sem modificação alguma.

Esse desenho tem uma limitação estrutural, no entanto, que motivou
diretamente o trabalho aqui descrito. O caminho ponta a ponta da rota
original é:

\begin{center}
malha nativa global $\to$ \textbf{grade lat-lon regular} $\to$ formato
binário WPS $\to$ malha nativa regional
\end{center}

Duas etapas desse caminho envolvem reamostragem espacial redundante: a
malha nativa global é primeiro remapeada para uma grade lat-lon regular
(pelo utilitário \texttt{convert\_mpas}, da NCAR), e depois o
\texttt{init\_atmosphere\_model} reinterpola essa grade de volta para a
malha nativa \emph{regional} de destino, usando interpolação bilinear ou
de 16 pontos sobre índices de grade estruturada -- o único modo de
interpolação horizontal que esse programa aceita para o caso de dado
real (\texttt{config\_init\_case=7/9}). O resultado é um percurso
\emph{malha $\to$ grade $\to$ malha}: cada etapa de reamostragem introduz
um erro de interpolação de ida e volta, mesmo quando a grade
intermediária é fina, simplesmente pelo fato de trocar de representação
espacial duas vezes em vez de zero. Esse problema já havia se manifestado
concretamente no projeto como uma classe de bug real e documentada (a
grade lat-lon nominal, dimensionada pela elipse teórica do recorte
regional, ficava sistematicamente menor que a extensão real da malha
recortada, deixando células sem dado válido na fronteira -- ver
\cref{cap:discussao}).

\section{Pergunta de pesquisa e escopo do trabalho}

A pergunta que orienta este trabalho é direta: \emph{é possível eliminar
inteiramente a grade lat-lon intermediária, interpolando diretamente
entre a malha de Voronoi de origem e a malha de Voronoi de destino, e
escrevendo os arquivos de entrada do núcleo dinâmico do \mpasa{}
(\texttt{init.nc}/\texttt{lbc.*.nc}) diretamente -- sem passar pelo
formato binário do WPS nem pelo \texttt{init\_atmosphere\_model}
original?}

A investigação começou pela leitura do código-fonte real do
\texttt{init\_atmosphere\_model} (rotina
\texttt{mpas\_init\_atm\_hinterp.F}) para verificar se essa reinterpolação
final para a malha regional aceitaria, de alguma forma, uma lista de
pontos dispersos em vez de uma grade estruturada. A resposta encontrada
foi negativa: o programa é estruturalmente amarrado ao formato de grade
(a interpolação bilinear opera sobre índices \texttt{(i,j)} de uma matriz
regular). Isso significa que \emph{qualquer} ganho de precisão de uma
interpolação nativa malha-a-malha só é alcançável pulando esse programa
inteiramente para a etapa de horizontal/vertical/balanço físico --- não
há um ``escopo mínimo'' que preserve o formato WPS e ainda elimine o
percurso malha$\to$grade$\to$malha. Essa constatação (detalhada na
\cref{cap:arquitetura}) definiu o escopo do trabalho: reimplementar, em
Fortran, cada etapa fisicamente necessária para produzir um
\texttt{init.nc}/\texttt{lbc.*.nc} válido -- não um subconjunto delas.

\section{Contribuições}

Este trabalho contribui com:

\begin{enumerate}[label=(\roman*)]
  \item Um motor de interpolação horizontal nativa malha-a-malha, que
    reaproveita sem modificação o mecanismo de pesos baricêntricos do
    utilitário \texttt{convert\_mpas} (originalmente exposto apenas para
    remapear para grades lat-lon), operando em modo de lista de pontos
    dispersos (\emph{scatter}) sobre os centros de célula da malha de
    destino;
  \item A extração e tradução literal, para um conjunto de módulos
    Fortran independentes e testáveis, das rotinas do núcleo dinâmico do
    MPAS-Model (versão de referência 3.0.2) responsáveis pela geração da
    grade vertical \textit{terrain-following} híbrida, interpolação
    vertical, balanço hidrostático não linear, e inicialização de campos
    de superfície e solo;
  \item Um escritor de \texttt{init.nc}/\texttt{lbc.*.nc} que produz
    arquivos NetCDF binariamente compatíveis com o que o
    \texttt{mpas\_atmosphere} espera, sem depender do
    \texttt{init\_atmosphere\_model} original para gerá-los;
  \item Uma generalização dos programas resultantes para ler os
    parâmetros de configuração (\texttt{config\_*}) do
    \texttt{namelist.init\_atmosphere} real do experimento em tempo de
    execução, tornando a rota reprodutível para qualquer malha e caso de
    estudo, não apenas para o exemplo de validação usado neste relatório;
  \item Validação em duas camadas: (a) numérica, campo a campo, contra
    arquivos \texttt{init.nc}/\texttt{lbc.*.nc} reais de produção; e
    (b) funcional, executando o \texttt{mpas\_atmosphere} real a partir
    dos arquivos gerados por essa rota e obtendo uma previsão de 24 horas
    fisicamente sã, sem erros.
\end{enumerate}

\section{Organização do relatório}

O \cref{cap:fundamentacao} apresenta a fundamentação teórica necessária
-- a estrutura da malha SCVT do \mpasa{}, sua dualidade com a
triangulação de Delaunay, e o princípio da interpolação baricêntrica
usada como método de remapeamento horizontal. O \cref{cap:literatura}
situa esse método na literatura mais ampla sobre malhas de Voronoi em
geociências e computação científica. O \cref{cap:arquitetura} descreve a
arquitetura geral da nova rota (o \emph{pipeline}) e a decisão de
implementação de escopo completo. Os capítulos~\ref{cap:fase1} a
\ref{cap:fase7} descrevem cada fase da implementação em detalhe, com as
equações e algoritmos efetivamente usados. O \cref{cap:implementacao}
discute as escolhas de engenharia da implementação em Fortran. O
\cref{cap:resultados} apresenta os resultados de validação numérica e a
execução funcional do modelo. O \cref{cap:discussao} discute limitações
conhecidas e simplificações deliberadas. O \cref{cap:conclusao} encerra
com as conclusões e sugestões de trabalho futuro.

\chapter{Fundamentação teórica}
\label{cap:fundamentacao}

\section{A malha do MPAS-A: tesselação de Voronoi centroidal esférica}

A malha horizontal do \mpasa{} é uma \emph{tesselação de Voronoi
centroidal esférica} (\textit{Spherical Centroidal Voronoi Tessellation},
SCVT) \cite{skamarock2012voronoi, dugunzburger2002cvt, yang2017cvtsphere}.
Dado um conjunto de pontos geradores $\{\mathbf{x}_i\}_{i=1}^{N}$
distribuídos sobre a esfera, o diagrama de Voronoi particiona a
superfície em $N$ células $\{V_i\}$, onde cada célula é o lugar
geométrico dos pontos da esfera mais próximos do gerador $\mathbf{x}_i$
do que de qualquer outro gerador:
\begin{equation}
  V_i = \left\{ \mathbf{x} \in \mathbb{S}^2 \;:\; d(\mathbf{x}, \mathbf{x}_i) \le d(\mathbf{x}, \mathbf{x}_j) \ \forall j \ne i \right\},
  \label{eq:voronoi_def}
\end{equation}
com $d(\cdot,\cdot)$ a distância geodésica sobre a esfera. A tesselação é
dita \emph{centroidal} quando cada gerador $\mathbf{x}_i$ coincide com o
centroide (baricentro de massa) de sua própria célula $V_i$ -- uma
condição de otimalidade obtida iterativamente pelo algoritmo de Lloyd, e
que produz células com forma próxima de hexágonos regulares na maior
parte do domínio, minimizando distorção de área e permitindo refinamento
local suave da resolução horizontal sem as descontinuidades abruptas de
grades aninhadas tradicionais.

No \mpasa{}, as variáveis prognósticas de massa (densidade seca $\rho$,
temperatura potencial $\theta$, razões de mistura de água $q$) são
armazenadas nos centros de célula $\mathbf{x}_i$ (geradores da
tesselação), enquanto a componente \emph{normal} do vento horizontal $u$
é armazenada nas arestas entre células vizinhas -- um posicionamento tipo
\emph{C-grid} (Arakawa C) generalizado para malha não estruturada. Essa
convenção de estagiamento é o que motiva a separação, na interpolação
horizontal nativa apresentada neste trabalho (\cref{cap:fase1}), entre um
tratamento para campos de massa (posicionados no centro da célula) e um
tratamento posterior, distinto, para a reconstrução do vento nas arestas
(\cref{cap:fase3}).

\section{Dualidade com a triangulação de Delaunay}
\label{sec:dualidade}

Toda tesselação de Voronoi possui uma estrutura geométrica \emph{dual}: a
\emph{triangulação de Delaunay}, obtida conectando por um segmento (ou,
na esfera, um arco geodésico) cada par de geradores cujas células de
Voronoi compartilham uma aresta. A propriedade central dessa dualidade,
usada diretamente neste trabalho, é que os \emph{vértices} da malha de
Voronoi (pontos onde três ou mais células se encontram) são exatamente os
\emph{circuncentros} dos triângulos da triangulação de Delaunay --
enquanto os \emph{centros das células} de Voronoi (os geradores
$\mathbf{x}_i$) são exatamente os \emph{vértices} dos triângulos de
Delaunay.

\begin{figure}[H]
  \centering
  \begin{tikzpicture}[scale=1.15, every node/.style={font=\small}]
    % --- celulas de Voronoi (poligonos irregulares esquematicos) ---
    \begin{scope}
      \clip (-4.6,-2.6) rectangle (4.6,2.6);
      % geradores (centros de celula)
      \coordinate (A) at (-1.4, 0.9);
      \coordinate (B) at ( 1.3, 1.1);
      \coordinate (C) at ( 0.1,-1.3);
      \coordinate (D) at (-2.6,-1.1);
      \coordinate (E) at ( 2.8,-0.6);
      \coordinate (F) at (-0.3, 2.4);
      \coordinate (G) at (-3.0, 1.6);
      \coordinate (Hp) at ( 3.2, 1.9);
      \coordinate (I) at ( 1.4,-2.5);
      \coordinate (J) at (-1.9,-2.6);

      % arestas de Voronoi (esquematicas, desenhadas a mao para ilustrar hexagonos)
      \draw[thick, blue!60!black]
        (-0.05,2.05) -- (1.55,2.55)
        (1.55,2.55) -- (2.55,1.15)
        (2.55,1.15) -- (2.05,-0.35)
        (2.05,-0.35) -- (0.75,-0.55)
        (0.75,-0.55) -- (-0.05,0.55)
        (-0.05,0.55) -- (-0.05,2.05);
      \draw[thick, blue!60!black]
        (-0.05,0.55) -- (-1.55,0.05)
        (-1.55,0.05) -- (-1.65,-1.35)
        (-1.65,-1.35) -- (-0.55,-1.95)
        (-0.55,-1.95) -- (0.75,-0.55);
      \draw[thick, blue!60!black]
        (-1.55,0.05) -- (-2.35, 1.35)
        (-2.35,1.35) -- (-0.05,2.05);
      \draw[thick, blue!60!black]
        (0.75,-0.55) -- (1.55,-1.75)
        (1.55,-1.75) -- (2.05,-0.35);
      \draw[thick, blue!60!black]
        (-1.65,-1.35) -- (-2.95,-0.85);
      \draw[thick, blue!60!black]
        (-0.55,-1.95) -- (-0.35,-2.95);
      \draw[thick, blue!60!black]
        (1.55,-1.75) -- (1.35,-2.9);
      \draw[thick, blue!60!black]
        (1.55,2.55) -- (2.95,2.45);
      \draw[thick, blue!60!black]
        (2.55,1.15) -- (3.55,1.55);
      \draw[thick, blue!60!black]
        (-2.35,1.35) -- (-3.55,1.75);
    \end{scope}

    % --- geradores / centros de celula ---
    \foreach \p/\lbl in {A/$\mathbf{x}_A$, B/$\mathbf{x}_B$, C/$\mathbf{x}_C$}
      \fill[red!70!black] (\p) circle (2.2pt);
    \foreach \p in {D,E,F,G,Hp,I,J}
      \fill[red!70!black] (\p) circle (1.6pt);

    % --- triangulo dual de Delaunay em torno do ponto de interesse ---
    \coordinate (P) at (0.05, 0.35);
    \draw[thick, dashed, orange!80!black] (A) -- (B) -- (C) -- cycle;
    \fill[orange!15] (A) -- (B) -- (C) -- cycle;
    \draw[thick, dashed, orange!80!black] (A) -- (B) -- (C) -- cycle;

    % ponto de interesse (ex.: centro de celula da malha de DESTINO)
    \fill[black] (P) circle (2.4pt);
    \node[below right=1pt] at (P) {$\mathbf{p}$ (destino)};

    \node[above=2pt] at (A) {$\mathbf{x}_A$};
    \node[above right=1pt] at (B) {$\mathbf{x}_B$};
    \node[below=2pt] at (C) {$\mathbf{x}_C$};

    \node[align=left, anchor=west, font=\footnotesize] at (3.7,-1.7)
      {\textcolor{blue!60!black}{\rule{1em}{1pt}} malha primal (Voronoi)\\[2pt]
       \textcolor{orange!80!black}{\rule{1em}{1pt}} malha dual (Delaunay)\\[2pt]
       \textcolor{red!70!black}{\textbullet} centro de célula};
  \end{tikzpicture}
  \caption{Esquema da dualidade Voronoi--Delaunay usada na interpolação
    horizontal nativa. As células poligonais (azul) formam a malha
    primal de Voronoi; conectando os centros de três células mutuamente
    vizinhas obtém-se um triângulo do dual de Delaunay (laranja). Um
    ponto arbitrário $\mathbf{p}$ -- tipicamente o centro de célula da
    malha de \emph{destino} -- é localizado dentro de um único triângulo
    de Delaunay da malha de \emph{origem}, e o campo escalar nesse ponto
    é obtido por interpolação baricêntrica dos três valores nos vértices
    do triângulo (\cref{eq:baricentrica}).}
  \label{fig:voronoi_delaunay}
\end{figure}

\section{Interpolação baricêntrica no triângulo dual}
\label{sec:baricentrica}

Dado um triângulo de Delaunay com vértices $\mathbf{x}_A, \mathbf{x}_B,
\mathbf{x}_C$ (centros de três células de Voronoi mutuamente vizinhas) e
um ponto de interesse $\mathbf{p}$ localizado dentro desse triângulo, as
\emph{coordenadas baricêntricas} $(\lambda_A, \lambda_B, \lambda_C)$ de
$\mathbf{p}$ são os únicos pesos que satisfazem simultaneamente
\begin{equation}
  \mathbf{p} = \lambda_A \mathbf{x}_A + \lambda_B \mathbf{x}_B + \lambda_C \mathbf{x}_C,
  \qquad
  \lambda_A + \lambda_B + \lambda_C = 1,
  \qquad \lambda_A, \lambda_B, \lambda_C \ge 0.
  \label{eq:baricentrica_def}
\end{equation}
Geometricamente, cada peso $\lambda_A$ é proporcional à área do
subtriângulo oposto ao vértice $A$ (formado por $\mathbf{p}$, $B$ e $C$),
normalizada pela área total do triângulo $ABC$:
\begin{equation}
  \lambda_A = \frac{\text{Área}(\mathbf{p}, B, C)}{\text{Área}(A,B,C)},
  \qquad
  \lambda_B = \frac{\text{Área}(A, \mathbf{p}, C)}{\text{Área}(A,B,C)},
  \qquad
  \lambda_C = \frac{\text{Área}(A, B, \mathbf{p})}{\text{Área}(A,B,C)}.
  \label{eq:baricentrica_area}
\end{equation}
O valor interpolado de um campo escalar $\phi$ (por exemplo, temperatura,
altura geopotencial, ou umidade específica) no ponto $\mathbf{p}$ é então
a combinação linear
\begin{equation}
  \phi(\mathbf{p}) \approx \lambda_A\, \phi(\mathbf{x}_A) + \lambda_B\, \phi(\mathbf{x}_B) + \lambda_C\, \phi(\mathbf{x}_C).
  \label{eq:baricentrica}
\end{equation}
Esse método é exato para campos que variam linearmente sobre o triângulo
e, criticamente para este trabalho, é \emph{exato por construção} quando
o ponto de interesse $\mathbf{p}$ coincide com um dos próprios vértices
$\mathbf{x}_A$, $\mathbf{x}_B$ ou $\mathbf{x}_C$ (nesse caso um dos pesos
vale $1$ e os outros dois valem $0$) -- é essa propriedade que permite
validar a implementação de forma direta: interpolando o campo de uma
malha para si mesma (\emph{auto-\emph{scatter}}), o erro deve ser
identicamente nulo em todas as células (\cref{sec:validacao_fase1}).

Na esfera, o procedimento operacional usado pelo motor de pesos (herdado
sem modificação do utilitário \texttt{convert\_mpas} da NCAR, ver
\cref{sec:fase1_reuso}) é: (i)~localizar, para cada ponto de destino
$\mathbf{p}$, o triângulo de Delaunay da malha de origem que o contém,
usando a conectividade célula-célula/aresta já armazenada na malha
(\texttt{cellsOnCell}, \texttt{cellsOnVertex}); (ii)~calcular os três
pesos baricêntricos via \cref{eq:baricentrica_area}, com as áreas
computadas sobre a superfície esférica; e (iii)~aplicar
\cref{eq:baricentrica} a cada campo escalar requisitado. A reconstrução
do vento, por armazenar sua componente normal nas arestas em vez de um
vetor completo no centro de célula, exige um tratamento adicional --
descrito em detalhe na \cref{sec:vento_arestas} -- baseado em Funções de
Base Radial (RBF) para reconstrução vetorial nos centros de célula,
seguida de projeção na direção normal de cada aresta de destino.

\section{Coordenada vertical \textit{terrain-following} híbrida}
\label{sec:coord_vertical_teoria}

Na direção vertical, o \mpasa{} usa uma coordenada $\zeta$
\emph{terrain-following} (segue o relevo perto da superfície e tende a
superfícies de altura constante longe dele), na formulação híbrida de
Klemp \cite{klemp2011terrainfollowing}. A altura física $z$ de uma
superfície de coordenada constante $\zeta=\zeta_k$ é dada por
\begin{equation}
  z(\zeta_k, \mathbf{x}) = \zeta_k + A(\zeta_k)\, h_s(\mathbf{x}, \zeta_k),
  \label{eq:coord_hibrida_teoria}
\end{equation}
onde $h_s$ é uma versão \emph{suavizada} do terreno (progressivamente
mais suave em níveis mais altos, para atenuar erros de gradiente de
pressão horizontal sobre relevo íngreme -- o problema clássico que
motivou toda a família de coordenadas híbridas desde Gal-Chen e
Simmons/Burridge até Klemp) e $A(\zeta)$ é uma função de atenuação que
vale $1$ na superfície e decai suavemente a $0$ numa altura de transição
$z_t$, definindo o ponto acima do qual as superfícies de coordenada
voltam a ser perfeitamente horizontais. As formas específicas de
$A(\zeta)$ e do algoritmo de suavização de $h_s$ efetivamente usadas pelo
\mpasa{} -- e reproduzidas neste trabalho -- são detalhadas com suas
constantes exatas no \cref{cap:fase2}.

\section{Balanço hidrostático de referência}
\label{sec:hidrostatico_teoria}

A condição inicial de um modelo não hidrostático como o \mpasa{} ainda
precisa de um estado \emph{base} hidrostaticamente balanceado, em torno
do qual as variáveis prognósticas são definidas como perturbações. O
\mpasa{} usa uma atmosfera de referência isotérmica seca
($T_0=\SI{250}{\kelvin}$), e o campo de pressão real é obtido resolvendo
iterativamente a relação hidrostática discreta na direção vertical --
não uma integração fechada, mas um esquema numérico que usa exatamente
os mesmos coeficientes de interpolação vertical ($f_{zm}$, $f_{zp}$) já
definidos pela grade vertical da \cref{sec:coord_vertical_teoria}, de
modo que o campo de massa final seja consistente, ponto a ponto, com a
métrica vertical adotada. O algoritmo completo, incluindo a formulação
não trivial de densidade seca acoplada à métrica vertical
($\rho_{zz}=\rho/\partial z/\partial \zeta$), está detalhado no
\cref{cap:fase4}.

\chapter{Revisão da literatura}
\label{cap:literatura}

Este capítulo situa o método adotado dentro da literatura mais ampla
sobre geração e uso de malhas de Voronoi, organizada em três grupos
temáticos: (i)~o núcleo dinâmico do próprio \mpasa{} e seu uso em
assimilação de dados; (ii)~geração e otimização de malhas de Voronoi
esféricas; e (iii)~interpolação sobre malhas de Voronoi/Delaunay fora do
domínio da geociência atmosférica, que ajudam a generalizar e comparar o
método baricêntrico adotado.

\section{O núcleo dinâmico MPAS-A e a malha SCVT}

A descrição de referência do núcleo dinâmico do \mpasa{}
\cite{skamarock2012voronoi} estabelece a malha $C$-\textit{grid} sobre
tesselação de Voronoi centroidal esférica como escolha estrutural do
modelo, motivada pela ausência de singularidades polares e pela
possibilidade de refinamento local suave -- exatamente as duas
propriedades que tornam a malha adequada tanto para a simulação quanto,
como demonstrado neste trabalho, para servir como origem/destino de uma
interpolação nativa sem grade intermediária. A apresentação de Duda
\cite{duda2014scalability} complementa essa descrição com considerações
de desempenho/escalabilidade da implementação de produção, relevantes
para justificar por que operações vetorizadas sobre a conectividade da
malha (como o cálculo de pesos baricêntricos célula a célula) são
tratáveis computacionalmente mesmo em malhas de centenas de milhares de
células.

O uso de interpolação baricêntrica na malha dual de Delaunay como método
de remapeamento entre estados do \mpasa{} já é prática estabelecida na
comunidade -- não uma proposta original deste trabalho, mas a aplicação
de um método já validado a um novo contexto (geração de condição
inicial/fronteira lateral, em vez de assimilação de dados). A referência
central para essa confirmação é Ha et al.
\cite{ha2017enkf}, que descreve o sistema de assimilação por
filtro de Kalman em conjunto (MPAS-DART) do NCAR: a Eq.~6 e a Fig.~1
daquele trabalho descrevem exatamente a interpolação baricêntrica no
triângulo dual de Delaunay (três centros de célula vizinhos) para campos
escalares de massa, e a Fig.~5 demonstra que a reconstrução vetorial do
vento nos centros de célula via Funções de Base Radial (RBF),
\emph{antes} de interpolar, produz resultado superior a tentar
interpolar a componente normal diretamente nas arestas -- é essa
constatação que justifica a ordem de operações adotada na
\cref{sec:vento_arestas} deste relatório. Como contraponto, os materiais
técnicos do sistema de assimilação MPAS-JEDI \cite{mpasjedi2021staticb}
revelam que mesmo abordagens mais recentes de assimilação variacional
ainda recorrem, para operações específicas sem inversa natural na malha
não estruturada (como a transformação $\psi,\chi \to u,v$ via
transformada espectral), a uma etapa intermediária em grade
latitude-longitude -- evidenciando que ``100\% nativo'' nem sempre é
possível para \emph{toda} operação de um sistema de modelagem/assimilação
completo, mas é possível, e é o padrão estabelecido, especificamente para
o remapeamento direto de campos entre duas malhas -- o problema atacado
neste trabalho.

\section{Geração e otimização de malhas de Voronoi esféricas}

A construção computacional de tesselações de Voronoi centroidais é, por
si, uma área de pesquisa ativa. Du e Gunzburger
\cite{dugunzburger2002cvt} formalizam a otimização de CVTs via
minimização de uma função energia relacionada ao algoritmo de Lloyd,
fundamento teórico do processo de geração da própria malha do \mpasa{}.
Yang, Gunzburger e Ju \cite{yang2017cvtsphere} propõem um método de
geração paralela mais rápido, baseado em um algoritmo LBFGS
pré-condicionado por Lloyd, especificamente para o caso esférico --
relevante para a geração de malhas regionais de alta resolução, embora
não usado diretamente na etapa de \emph{interpolação} deste trabalho (que
opera sobre malhas já geradas por outra ferramenta,
\texttt{MPAS-Limited-Area}). Engwirda \cite{engwirda2017jigsawgeo}
descreve o \textsc{jigsaw-geo}, um gerador de malha não estruturada
ortogonal localmente escalonada para modelagem de circulação geral,
alternativa de geração de malha usada por outros modelos da mesma
família dinâmica (por exemplo MPAS-Ocean); e Santos e Peixoto
\cite{santospeixoto2021refinement} exploram refinamento local de malha
esférica de Voronoi orientado por topografia em modelos de água rasa,
tema adjacente ao problema de suavização de terreno tratado na
\cref{cap:fase2}\footnote{Refinamento de \emph{malha} (mudar a resolução
espacial dos geradores conforme o relevo) e suavização do \emph{campo de
terreno} usado na definição da coordenada vertical (mantendo a malha
fixa, o problema efetivamente tratado neste trabalho) são technicamente
distintos, mas motivados pelo mesmo fenômeno físico -- erro de gradiente
de pressão horizontal sobre relevo íngreme.}.

\section{Interpolação sobre malhas de Voronoi/Delaunay em outros domínios}

Fora da modelagem atmosférica, a interpolação sobre malhas de Voronoi tem
tradição estabelecida em geometria computacional e em outras áreas da
física computacional, o que permite contextualizar o método baricêntrico
como caso particular de uma família mais ampla de interpolantes.
Hiyoshi e Sugihara \cite{hiyoshi2000continuity} apresentam um arcabouço
geral de interpolantes baseados no diagrama de Voronoi com continuidade
$C^1$ (generalizando os interpolantes clássicos de Sibson e de Laplace),
mais suaves que a interpolação baricêntrica linear por partes usada
neste trabalho -- um candidato natural de evolução futura caso a
descontinuidade de derivada nas arestas dos triângulos de Delaunay se
mostre relevante em aplicações futuras (\cref{cap:discussao}). Melodia e
Lenz \cite{melodia2019persistenthomology} propõem um critério de parada
baseado em homologia persistente para refinar iterativamente uma
interpolação de Voronoi, tema de otimização de malha não diretamente
aplicável ao caso deste trabalho (malhas fixas, pré-geradas), mas
ilustrativo da maturidade do campo. Em astrofísica e ciências da Terra,
Petkova et al.~\cite{petkova2017densitygridding} usam grades de Voronoi
para regularizar dados de hidrodinâmica lagrangiana (SPH), e Camps et
al.~\cite{camps2013radiativetransfer} e Bonduà et al.
\cite{bondua2017voronoi3d} aplicam malhas de Voronoi tridimensionais em
transferência radiativa astrofísica e migração de gás em formações
sedimentares, respectivamente -- ambos os casos usando a mesma
propriedade central explorada neste trabalho (adaptabilidade local de
resolução sem descontinuidade de malha), embora em três dimensões
espaciais em vez da superfície esférica bidimensional do \mpasa{}.
Abdelkader et al.~\cite{abdelkader2019vorocrust} tratam do problema
correlato, mas distinto, de geração de malha de Voronoi conformada a uma
superfície sem recorte (\emph{clipping}) dos poliedros de fronteira --
não diretamente aplicável aqui, já que a malha do \mpasa{} já vem
pré-gerada e recortada pela ferramenta \texttt{MPAS-Limited-Area}, mas
relevante como pano de fundo sobre os desafios numéricos de se trabalhar
com malhas de Voronoi em domínios limitados.

\section{Síntese}

A literatura revisada confirma dois pontos que sustentam as decisões de
projeto deste trabalho: primeiro, que a interpolação baricêntrica na
malha dual de Delaunay não é uma escolha \emph{ad hoc}, mas o método
padrão já validado pela comunidade MPAS (via MPAS-DART) para
remapeamento nativo de campos de massa entre estados dessa mesma malha;
segundo, que existem interpolantes de maior continuidade
(Hiyoshi--Sugihara) e técnicas mais sofisticadas de geração/refinamento
de malha que não foram necessários para o escopo deste trabalho -- que
lida com remapeamento entre duas malhas \emph{já geradas} -- mas que
constituem direções concretas de trabalho futuro caso a precisão do
método baricêntrico linear se mostre insuficiente em aplicações que
exijam maior suavidade de derivadas (por exemplo, esquemas de advecção
de ordem mais alta que dependam de gradientes horizontais contínuos do
campo interpolado).

\chapter{Arquitetura da rota de interpolação nativa}
\label{cap:arquitetura}

\section{A rota original e sua limitação estrutural}
\label{sec:rota_original}

O \selfgrib{} original gera condição inicial e fronteira lateral em três
estágios, ilustrados na metade superior da \cref{fig:comparacao_rotas}:
\texttt{extract\_fields} interpola verticalmente os campos nativos do
\mpasa{} (55 níveis de modelo) para um conjunto de 61 níveis de pressão
fixa, ainda sobre a malha nativa \emph{global}; \texttt{convert\_mpas}
(ferramenta da NCAR, reaproveitada sem modificação) remapeia
horizontalmente esses campos da malha nativa para uma grade
latitude-longitude regular; e \texttt{pack\_intermediate} escreve o
resultado no formato binário intermediário do WPS. Esse arquivo
intermediário é então consumido pelo \texttt{init\_atmosphere\_model}
original (inalterado), que interpola de volta -- bilinear ou 16 pontos,
sobre índices de grade estruturada -- para a malha nativa \emph{regional}
de destino, e resolve internamente a grade vertical, a interpolação
vertical e o balanço hidrostático.

Como discutido na \cref{cap:introducao}, essa rota tem uma limitação
estrutural: o \texttt{init\_atmosphere\_model} não aceita, para o caso de
dado real, nenhuma forma de entrada que não seja uma grade estruturada
-- confirmado pela leitura da rotina \texttt{mpas\_init\_atm\_hinterp.F}
do código-fonte de referência. Isso significa que o ganho de precisão de
uma interpolação nativa malha-a-malha (\cref{sec:baricentrica}) só é
alcançável substituindo \emph{inteiramente} o caminho
\texttt{convert\_mpas}~$\to$~formato WPS~$\to$~\texttt{init\_atmosphere\_model},
não apenas uma parte dele.

\section{A rota nativa: visão geral}
\label{sec:rota_nativa}

A rota descrita neste relatório (metade inferior da
\cref{fig:comparacao_rotas}) substitui esse caminho por sete fases,
implementadas como programas Fortran independentes que se comunicam por
arquivos NetCDF intermediários -- mantendo a mesma filosofia de
composição em estágios discretos e idempotentes já usada na rota
original, mas eliminando toda e qualquer grade regular do caminho de
dados.

\begin{figure}[H]
  \centering
  \begin{tikzpicture}[
      node distance=6mm and 10mm,
      box/.style={draw, rounded corners, minimum height=0.9cm, minimum width=2.6cm, align=center, font=\footnotesize, fill=blue!6},
      grid/.style={draw, minimum height=0.9cm, minimum width=2.6cm, align=center, font=\footnotesize, fill=red!8, dashed},
      lbl/.style={font=\scriptsize\itshape, text width=2.6cm, align=center},
      arr/.style={-{Stealth[length=2mm]}, thick},
    ]

    % ===================== ROTA ORIGINAL =====================
    \node[font=\bfseries\small] (titulo1) at (0,4.6) {Rota original (via WPS)};

    \node[box] (g1) at (-6.5,3.6) {malha nativa\\GLOBAL};
    \node[box] (e1) at (-3.7,3.6) {\texttt{extract\_}\\\texttt{fields}};
    \node[grid] (c1) at (-0.9,3.6) {grade\\lat-lon};
    \node[box] (p1) at (1.9,3.6)  {formato\\WPS};
    \node[box] (i1) at (4.9,3.6)  {\texttt{init\_atmosphere\_}\\\texttt{model} (original)};
    \node[box] (r1) at (7.9,3.6)  {malha nativa\\REGIONAL};

    \draw[arr] (g1) -- (e1);
    \draw[arr] (e1) -- (c1) node[midway, above, font=\tiny] {\texttt{convert\_mpas}};
    \draw[arr] (c1) -- (p1) node[midway, above, font=\tiny] {\texttt{pack\_intermediate}};
    \draw[arr] (p1) -- (i1);
    \draw[arr] (i1) -- (r1) node[midway, above, font=\tiny] {bilinear/16-pt};

    \node[lbl, red!50!black] at (-0.9,2.75) {\textbf{grade intermediária}\\(fonte do erro de\\ida-e-volta)};

    % ===================== ROTA NATIVA =====================
    \node[font=\bfseries\small] (titulo2) at (0,1.4) {Rota nativa (\emph{voronoi-to-voronoi}, este trabalho)};

    \node[box] (g2) at (-6.5,0.4) {malha nativa\\GLOBAL};
    \node[box, fill=green!10] (f1) at (-3.7,0.4) {Fase 1\\\texttt{hinterp\_native}};
    \node[box, fill=green!10] (f2) at (-0.9,0.4) {Fase 2\\grade vertical};
    \node[box, fill=green!10] (f34) at (1.9,0.4) {Fases 3+4\\interp. vert. +\\hidrostático};
    \node[box, fill=green!10] (f6) at (4.9,0.4) {Fase 6\\superfície/solo};
    \node[box] (r2) at (7.9,0.4) {\texttt{init.nc} /\\\texttt{lbc.*.nc}};

    \draw[arr] (g2) -- (f1);
    \draw[arr] (f1) -- (f2);
    \draw[arr] (f2) -- (f34);
    \draw[arr] (f34) -- (f6);
    \draw[arr] (f6) -- (r2);

    \node[box, fill=green!10] (f7) at (7.9,-0.9) {Fase 7\\\texttt{mpas\_atmosphere}};
    \draw[arr] (r2) -- (f7);

    \node[lbl, green!35!black] at (-3.7,-0.5) {interpolação\\baricêntrica\\direta malha--malha};

  \end{tikzpicture}
  \caption{Comparação entre a rota original (via formato binário WPS e
    grade lat-lon intermediária, acima) e a rota nativa
    \emph{voronoi-to-voronoi} apresentada neste trabalho (abaixo). A rota
    nativa elimina toda etapa de reamostragem em grade regular: da malha
    de Voronoi global de origem à malha de Voronoi regional de destino,
    o dado nunca deixa de estar sobre uma malha não estruturada.}
  \label{fig:comparacao_rotas}
\end{figure}

As sete fases, cujo detalhamento técnico ocupa os
Capítulos~\ref{cap:fase1}--\ref{cap:fase7}, são:

\begin{description}[leftmargin=2.2cm, style=nextline]
  \item[Fase 1 -- Interpolação horizontal nativa] Remapeia os campos já
    extraídos e verticalmente interpolados para níveis de pressão fixa
    (ainda na malha global) diretamente para os centros de célula da
    malha regional de destino, via interpolação baricêntrica
    (\cref{sec:baricentrica}), sem grade intermediária.
  \item[Fase 2 -- Grade vertical nativa] Calcula a geometria vertical
    \textit{terrain-following} híbrida da malha de destino
    ($z_{grid}$, $zz$, coeficientes de interpolação vertical, termos de
    métrica para advecção) a partir apenas do terreno e da conectividade
    dessa malha -- não depende do dado meteorológico.
  \item[Fase 3 -- Interpolação vertical] Interpola os campos já
    horizontalmente remapeados (ainda em níveis de pressão fixa) para os
    níveis de modelo da grade vertical nativa calculada na Fase~2.
  \item[Fase 4 -- Balanço hidrostático] Resolve o campo de pressão,
    densidade seca e temperatura potencial (úmida) de forma consistente
    com a métrica vertical, além de água precipitável e vento vertical
    diagnóstico.
  \item[Fase 6 -- Campos de superfície e solo] Calcula os campos que não
    vêm diretamente da coluna atmosférica: temperatura profunda do solo,
    correção de temperatura de pele por diferença de elevação, fração de
    vegetação/albedo de fundo por interpolação climatológica mensal,
    flag de neve/gelo marinho.
  \item[Fase 7 -- Escrita e execução] Escreve os arquivos
    \texttt{init.nc}/\texttt{lbc.*.nc} completos (mesclando os campos
    calculados nas fases anteriores com os campos de cópia direta da
    malha estática) e alimenta o núcleo \texttt{mpas\_atmosphere} real.
\end{description}

\begin{center}
\begin{minipage}{0.9\textwidth}
\small\itshape
Nota de numeração: a Fase~5 do plano de implementação original (projeção
de vento células$\to$arestas) foi absorvida dentro das Fases~3/4, já que
a reconstrução da componente normal do vento e sua multiplicação pela
densidade ($ru$) fazem parte, respectivamente, da interpolação vertical
e do balanço hidrostático -- não há um programa Fase~5 isolado.
\end{minipage}
\end{center}

\section{Decisão de escopo: por que não um subconjunto}
\label{sec:decisao_escopo}

Uma primeira hipótese de trabalho, descartada ainda na fase de
investigação, era um ``escopo mínimo'': gerar o dado meteorológico em
modo \emph{scatter} (lista de pontos exatos da malha regional, em vez de
grade), mas ainda escrever no formato binário WPS e deixar o
\texttt{init\_atmosphere\_model} original consumi-lo. A leitura do
código-fonte de \texttt{mpas\_init\_atm\_hinterp.F} mostrou que essa
hipótese está estruturalmente quebrada: o formato WPS \emph{é}, por
definição, uma grade regular (o próprio cabeçalho de cada registro do
formato carrega \texttt{nx}, \texttt{ny}, espaçamento e origem de uma
grade); não existe uma forma de ``passar direto'' pontos dispersos por
esse caminho. Mesmo alimentando o \texttt{convert\_mpas} em modo
\emph{scatter} para gerar um arquivo tecnicamente válido, o
\texttt{init\_atmosphere\_model} ainda executaria um segundo salto de
reinterpolação bilinear/16 pontos ao final -- o ganho de precisão
validado no protótipo inicial (recuperação exata nos centros de célula,
\cref{sec:validacao_fase1}) só é alcançável pulando esse caminho por
completo, não reduzindo seu escopo.

Por essa razão, a decisão de projeto foi pelo escopo completo: reescrever
cada etapa fisicamente necessária -- grade vertical, interpolação
vertical, balanço hidrostático, campos de superfície -- como módulos
Fortran novos, extraídos literalmente do código-fonte de referência do
MPAS-Model (\cref{cap:implementacao}), em vez de reaproveitar o
\texttt{init\_atmosphere\_model} para qualquer parte dessas etapas.

\section{Generalização: leitura do \textit{namelist} real em tempo de execução}
\label{sec:generalizacao}

Uma decisão de engenharia transversal a todas as fases, adotada após
revisão explícita do escopo do trabalho, foi banir do código-fonte
qualquer parâmetro de configuração (\texttt{config\_*}) fixo como
constante de compilação. Nas primeiras versões dos programas, valores
como \texttt{config\_ztop}, \texttt{config\_dzmin},
\texttt{config\_nsm}, e -- criticamente -- \texttt{config\_start\_time}
(a própria data/hora da rodada) estavam hard-coded, o que restringia a
validação a um único caso de teste e, mais grave, impedia diretamente a
geração da fronteira lateral (\texttt{lbc.*.nc}), já que cada arquivo de
fronteira corresponde a um tempo diferente dentro do mesmo experimento.
Um módulo dedicado (\texttt{namelist\_config.F90},
\cref{sec:namelist_config}) foi introduzido para ler o
\texttt{namelist.init\_atmosphere} \emph{real} do experimento, via
\texttt{NAMELIST} nativo do Fortran, em tempo de execução -- o mesmo
arquivo de configuração já usado pela rota original, sem introduzir
nenhum formato novo. Essa mudança, além de generalizar a ferramenta para
qualquer malha/experimento, revelou um bug real de suposição incorreta
sobre o valor padrão de \texttt{config\_frac\_seaice}
(\cref{cap:discussao}), que só se manifestaria em malhas com gelo
marinho -- não presente no caso de validação usado neste relatório, mas
corrigido preventivamente.

\chapter{Fase 1 -- Interpolação horizontal nativa}
\label{cap:fase1}

\section{Objetivo e reaproveitamento de código}
\label{sec:fase1_reuso}

O objetivo da Fase~1 é remapear, célula a célula, os campos escalares já
extraídos e verticalmente interpolados para níveis de pressão fixa
(estágio herdado sem modificação da rota original, programa
\texttt{extract\_fields}) da malha nativa \emph{global} de origem para os
centros de célula da malha nativa \emph{regional} de destino, aplicando
a interpolação baricêntrica descrita na \cref{sec:baricentrica}.

A decisão de engenharia central desta fase foi \emph{não reescrever} o
motor de cálculo de pesos baricêntricos e localização de ponto no
triângulo dual: o utilitário \texttt{convert\_mpas} (mantido como
dependência da rota original, propriedade da NCAR) já implementa esse
mecanismo internamente -- confirmado por leitura de código-fonte,
seção ``Interpolation methodology'' de sua documentação -- para o caso
de remapear \emph{para} uma grade latitude-longitude regular. A
investigação encontrou que o módulo interno responsável por montar a
malha de destino (\texttt{target\_mesh.F90}) já possui um modo
alternativo, não exposto pela interface de linha de comando original,
que aceita uma \emph{lista de pontos dispersos} (\emph{scatter}) em vez
de uma grade regular -- bastando fornecer os vetores de latitude e
longitude dos pontos de destino com \texttt{irank=1}. O motor de cálculo
de pesos propriamente dito (\texttt{remapper.F90}) é inteiramente
agnóstico entre os dois modos.

O programa \texttt{hinterp\_native.F90}, novo neste trabalho, é
essencialmente uma cópia da estrutura de \texttt{convert\_mpas.F90} (os
mesmos módulos de leitura de malha, varredura de campos e escrita de
saída), com uma única mudança: em vez de chamar a rotina de montagem de
grade regular, ele lê as coordenadas \texttt{latCell}/\texttt{lonCell}
da malha de destino e as passa ao módulo de montagem de malha em modo
\emph{scatter}. Nenhuma linha do motor de pesos baricêntricos
(\texttt{remapper.F90}) foi modificada.

\section{Procedimento}

Para cada campo escalar $\phi$ requisitado (temperatura, componentes de
vento, altura geopotencial, umidade específica e relativa, e as
respectivas variáveis de superfície), e para cada célula $\mathbf{p}_j$
da malha de destino:

\begin{enumerate}
  \item Localiza-se o triângulo de Delaunay da malha de origem que
    contém $\mathbf{p}_j$, usando a conectividade célula-célula/vértice
    já armazenada na malha de origem (o dual de Delaunay não precisa ser
    recalculado do zero -- é uma propriedade geométrica derivável
    diretamente de \texttt{cellsOnVertex}/\texttt{verticesOnCell}, já
    presentes no arquivo de malha).
  \item Calculam-se os três pesos baricêntricos $(\lambda_A, \lambda_B,
    \lambda_C)$ via \cref{eq:baricentrica_area}, com as áreas computadas
    sobre a superfície esférica.
  \item Aplica-se \cref{eq:baricentrica} para obter
    $\phi(\mathbf{p}_j)$.
\end{enumerate}

O resultado é escrito em um arquivo NetCDF intermediário
(\texttt{native\_target.nc}), com os campos ainda organizados por nível
de pressão fixa mas já localizados nos centros de célula da malha
regional -- pronto para a interpolação vertical da Fase~3
(\cref{cap:fase3}).

\begin{caixaachado}
Um efeito colateral do reúso do módulo \texttt{remapper.F90} sem
modificação foi que suas rotinas internas de geração das variáveis de
coordenada de saída (\texttt{remap\_get\_target\_latitudes}/
\texttt{longitudes}) assumem semântica de grade regular e, em modo
\emph{scatter}, colapsam a variável \texttt{latitude} de saída para um
único valor escalar em vez de um vetor por ponto -- um bug de uso do
módulo original (não corrigido no \texttt{convert\_mpas} em si) que
passou despercebido por não ser exercitado no caminho de código
tradicional (sempre grade). A correção adotada foi \emph{aditiva}, não
substitutiva: em vez de reescrever essas rotinas (uma primeira tentativa
nesse sentido quebrou a escrita do NetCDF, com erro de ordem de criação
de dimensões em modo \texttt{define}), o programa novo escreve campos
adicionais corretos, \texttt{target\_latitude}/\texttt{target\_longitude},
por cima da saída original -- preservando as chamadas herdadas intactas.
\end{caixaachado}

\section{Reconstrução do vento nas arestas}
\label{sec:vento_arestas}

O vento não é tratado como um campo escalar simples nesta fase, pelo
motivo estrutural já discutido na \cref{sec:dualidade}: sua componente
prognóstica no \mpasa{} é a componente \emph{normal} $u$ armazenada nas
arestas entre células, não um vetor $(u,v)$ no centro de célula. Seguindo
a evidência apresentada por Ha et al.~\cite{ha2017enkf} de que reconstruir
o vento a partir das componentes zonal/meridional originais do
first-guess (que já estão disponíveis, por definição, nos centros de
célula da malha de \emph{origem}) e só então projetar na direção normal
da aresta de \emph{destino} produz resultado superior a interpolar
diretamente uma quantidade já projetada, o procedimento adotado -- ainda
na etapa combinada com a Fase~3 (\cref{cap:fase3}) -- calcula, para cada
aresta $e$ da malha de destino com células vizinhas $c_1(e), c_2(e)$:
\begin{equation}
  u_{fg}(e) = \cos(\alpha_e)\, \bar{u}_{zonal} + \sin(\alpha_e)\, \bar{v}_{meridional},
  \qquad
  \bar{(\cdot)} = \tfrac{1}{2}\big[(\cdot)_{c_1(e)} + (\cdot)_{c_2(e)}\big],
  \label{eq:vento_projecao}
\end{equation}
onde $\alpha_e$ é o ângulo da aresta $e$ da malha de \emph{destino} em
relação ao meridiano local (variável \texttt{angleEdge}, já presente na
malha) e as médias $\bar{u}_{zonal}$, $\bar{v}_{meridional}$ são tomadas
sobre os valores \emph{já remapeados horizontalmente} (via
\cref{eq:baricentrica}, campo a campo) nos dois centros de célula
vizinhos da aresta de destino. Esse perfil vertical projetado é então
interpolado verticalmente (\cref{cap:fase3}) usando como coordenada de
altura a média das alturas geopotenciais das duas células vizinhas.

\section{Validação numérica}
\label{sec:validacao_fase1}

A validação desta fase explora diretamente a propriedade de exatidão nos
vértices discutida na \cref{sec:baricentrica}: interpolando os campos de
uma malha de teste (\texttt{x1.2562.static.nc}, 2562 células) para os
próprios centros de célula dessa mesma malha (\emph{auto-\emph{scatter}}),
o erro obtido foi \emph{identicamente nulo} em todas as 2562 células,
tanto nos valores de campo quanto nas coordenadas de saída
(\texttt{target\_latitude}/\texttt{target\_longitude}). Como comparação
de controle, o mesmo experimento repetido passando por uma grade lat-lon
intermediária fina ($\SI{0.25}{\degree}$) -- reproduzindo o caminho da
rota original -- introduziu um erro sistemático mensurável, da ordem de
$10^{-4}$ a $10^{-3}$ (em unidades do campo), \emph{mesmo nesses mesmos
pontos de coincidência exata}, confirmando de forma independente que o
erro de ida-e-volta discutido na \cref{cap:introducao} é real e
mensurável, não apenas uma preocupação teórica.

\chapter{Fase 2 -- Grade vertical nativa}
\label{cap:fase2}

\section{Objetivo}

A Fase~2 calcula a geometria vertical completa da malha de destino --
altura física das interfaces entre níveis de modelo ($z_{grid}$), a
métrica vertical $zz$, coeficientes de interpolação/derivação vertical, e
termos de métrica para advecção de terceira ordem -- a partir apenas do
terreno e da conectividade da malha de destino, \textbf{sem nenhuma
dependência do dado meteorológico}. Diferentemente das demais fases, essa
etapa é puramente geométrica e determinística, o que permitiu uma
estratégia de validação particularmente forte: comparar o resultado
diretamente contra a grade vertical de um \texttt{init.nc} real de
produção, sem depender de nenhuma rodada nova.

As fórmulas desta seção foram extraídas literalmente do bloco
\texttt{config\_tc\_vertical\_grid} da rotina \texttt{init\_atm\_case\_gfs}
do código-fonte de referência do MPAS-Model (versão de produção
\texttt{mpas-bundle-3.0.2}, confirmada como o fonte exato que gerou os
binários usados no cluster de produção via inspeção do
\texttt{CMakeCache.txt} do sistema de build) -- não uma reformulação, mas
uma tradução linha a linha, removendo apenas a infraestrutura de troca de
halo entre blocos MPI (desnecessária nesta implementação, que roda serial
sobre a malha inteira).

\section{Perfil vertical de referência $z_w(k)$}
\label{sec:zw}

O perfil de altura das interfaces de nível na configuração
``2014 TC experiments'' (identificada, por inspeção do \textit{namelist}
real de produção, como o branch de código efetivamente usado em todas as
rodadas deste projeto -- não a fórmula genérica mais simples de versões
anteriores do MPAS) é definido por trechos, com constantes calibradas
para $n_{VertLevels}=55$:
\begin{equation}
  z_w(k) =
  \begin{cases}
    \left[ a_s \dfrac{k-1}{n_1} + c_2\left(\dfrac{(k-1)\Delta z}{z_t \zeta_l}\right)^{\!2} - c_3\left(\dfrac{(k-1)\Delta z}{z_t \zeta_l}\right)^{\!3} \right] z_t, & \dfrac{k-1}{n_1} < \zeta_l \\[1.2em]
    \left[ z_l + a_t\left(\dfrac{k-1}{n_1} - \zeta_l\right) \right] z_t, & \dfrac{k-1}{n_1} \ge \zeta_l
  \end{cases}
  \label{eq:zw}
\end{equation}
com $n_1 = n_{VertLevels}$, $\Delta z = z_t/n_1$, $z_l = 1 - a_t(1-\zeta_l)$,
$c_2 = 3(1-a_t) + 2(a_t-a_s)\zeta_l$, $c_3 = 2(1-a_t) + (a_t-a_s)\zeta_l$,
e as constantes calibradas $a_s = 0{,}075$, $a_t = 1{,}70$,
$\zeta_l = 0{,}75$ (para $n_{VertLevels} \ge 55$; valores distintos,
$a_t=1{,}23$ e $\zeta_l=0{,}31$, aplicam-se ao caso legado de 41 níveis,
não usado neste projeto). O topo do domínio $z_t$ é o parâmetro de
configuração \texttt{config\_ztop} (\SI{30000}{\metre} em todos os casos
de produção deste projeto).

\section{Coordenada híbrida e o achado do bloqueio}
\label{sec:ah_hibrido}

A função de atenuação $A(\zeta)$ introduzida na
\cref{eq:coord_hibrida_teoria} determina a rapidez com que a influência
do terreno é removida das superfícies de coordenada com a altura. O
código-fonte de referência oferece dois ramos:
\begin{equation}
  A(k) =
  \begin{cases}
    \cos^6\!\left( \dfrac{\pi}{2} \dfrac{z_w(k)}{z_h} \right), & z_w(k) < z_h \quad \text{(coordenada híbrida)} \\[0.6em]
    0, & z_w(k) \ge z_h \\[0.6em]
    1 - \dfrac{z_w(k)}{z_t}, & \text{(coordenada \textit{terrain-following} básica)}
  \end{cases}
  \label{eq:ah}
\end{equation}
onde $z_h$ é a altura de transição (parâmetro \texttt{config\_hybrid\_top\_z}).

\begin{caixaachado}
\label{achado:hybrid}
Este ponto foi o principal bloqueio numérico enfrentado no
desenvolvimento desta fase, e vale ser documentado em detalhe pelo valor
metodológico da investigação. A suposição inicial, baseada na ausência
de qualquer menção a \texttt{config\_hybrid\_coordinate} nos arquivos
\texttt{namelist.init\_atmosphere} reais de produção, foi de que essa
opção usava seu valor padrão \texttt{.false.} -- isto é, o ramo
\textit{terrain-following} básico (segunda linha da
\cref{eq:ah}). Sob essa suposição, a grade vertical calculada
divergia da referência real em até $\sim\SI{2}{\kilo\metre}$ em níveis
intermediários -- um erro fisicamente inaceitável, da mesma classe de
gravidade dos bugs catastróficos já documentados no relatório principal
do projeto (NaN em \texttt{PSFC}).

Uma primeira hipótese, investigada extensivamente, foi de que o processo
de suavização iterativa (\cref{sec:suavizacao}) seria genuinamente não
reprodutível bit a bit entre uma reimplementação serial e a rodada real
distribuída por MPI, dada a natureza do critério de aceitação por limiar
usado nesse algoritmo (pequenas diferenças de arredondamento podem levar
células individuais a aceitar ou rejeitar a atualização de forma
diferente, compondo-se ao longo de dezenas de passadas). Essa hipótese
foi \emph{descartada} experimentalmente: reimplementar toda a cadeia de
suavização em precisão simples (\texttt{real*4}), para casar exatamente
com a precisão do binário de produção (confirmada pelo log real,
``Default real precision: single''), não produziu nenhuma melhora
mensurável -- os valores de erro permaneceram idênticos, dígito a dígito.

A causa raiz real só foi encontrada ao baixar o arquivo
\texttt{Registry.xml} do \texttt{core\_init\_atmosphere} diretamente do
repositório de referência do MPAS-Model: a opção
\texttt{config\_hybrid\_coordinate} tem \texttt{default\_value="true"}
(com \texttt{config\_hybrid\_top\_z} padrão \SI{30000}{\metre}) --
confirmado de forma ainda mais direta ao acessar o código-fonte exato
usado para compilar o binário de produção (\texttt{mpas-bundle-3.0.2}),
onde a coordenada híbrida está, na verdade, \emph{fixada no
código-fonte} (\texttt{hybrid = .true.}, não sequer exposta como opção de
\textit{namelist} nessa versão). Ou seja: a ausência da chave no
\textit{namelist} real não significava ``desativado'', significava
``usando o padrão -- que é ativado''. Corrigindo a fórmula de $A(k)$ para
o primeiro ramo da \cref{eq:ah}, o erro máximo na grade vertical caiu de
$\sim\SI{1103}{\metre}$ para $\SI{171}{\metre}$, e o erro médio de
$\SI{27,4}{\metre}$ para $\SI{0,054}{\metre}$ -- uma redução de
aproximadamente 500 vezes, com o resíduo restante concentrado quase
inteiramente no anel de células de fronteira da malha regional
(\cref{sec:validacao_fase2}), não mais distribuído por toda a coluna
vertical.

A lição metodológica: uma revisão linha a linha do bloco de código
suspeito (o laço de suavização) não encontrou nenhuma diferença de
lógica porque a causa real estava numa dependência \emph{upstream} desse
bloco (o cálculo de $A(k)$, calculado dezenas de linhas antes) -- ao
investigar uma divergência numérica que resiste a uma revisão de código
aparentemente correta, é necessário validar também as dependências do
trecho suspeito, não apenas o próprio trecho.
\end{caixaachado}

\section{Suavização do terreno e da grade vertical}
\label{sec:suavizacao}

Antes de aplicar as fórmulas anteriores, o campo de terreno bruto passa
por uma suavização espacial de quarta ordem (filtro de Shapiro,
$n_{smterrain}$ passadas, tipicamente $1$), com coeficiente $+0{,}216$
seguido de $-0{,}216$ em cada passada -- validado como essencialmente
exato ($\text{erro médio}\approx\SI{0,15}{\metre}$) contra a referência
real.

Mais crítico é o algoritmo de suavização \emph{iterativa por nível} da
superfície de terreno $h_x(k,\cdot)$, que progressivamente remove
estrutura de pequena escala do terreno conforme a altura aumenta -- a
técnica descrita originalmente por Klemp
\cite{klemp2011terrainfollowing} (Eqs.~5--9 daquele trabalho, confirmadas
como estruturalmente idênticas ao código real do MPAS-A nesta
investigação). Para cada nível $k=2,\dots,n_{VertLevels}-1$, executam-se
$n_{sm}+k$ passadas de um filtro Laplaciano local ponderado pela
resolução da malha,
\begin{equation}
  h_x^{(n)}(i) = h_x^{(n-1)}(i) + s(i)\!\!\sum_{j \in \mathcal{N}(i)}\!\! \frac{\ell_{ij}}{d_{ij}}\Big[h_x^{(n-1)}(j) - h_x^{(n-1)}(i)\Big],
  \label{eq:suavizacao_hx}
\end{equation}
onde $\mathcal{N}(i)$ é o conjunto de células vizinhas de $i$,
$\ell_{ij}$/$d_{ij}$ são o comprimento da aresta e a distância entre
centros de célula, e $s(i)$ é um fator de suavização local escalado pela
resolução nominal da malha e pela altura do nível corrente. A atualização
só é \emph{aceita} se o espaçamento vertical resultante não cair abaixo
de uma fração mínima configurável do espaçamento nominal
($\delta z_{min} > \texttt{config\_dzmin} \cdot \Delta z_w(k)$) --
critério que existe justamente para impedir camadas degeneradas
(espaçamento nulo ou negativo) perto de relevo íngreme.

\section{Campos derivados e coeficientes numéricos}

A partir de $z_w(k)$ e $A(k)$, calculam-se diretamente (sem depender de
mais nenhuma suavização): $\Delta z_w(k) = z_w(k{+}1)-z_w(k)$,
$rdz_w(k)=1/\Delta z_w(k)$, $\Delta z_u(k)=\tfrac12[\Delta z_w(k)+\Delta z_w(k{-}1)]$,
$rdz_u(k)=1/\Delta z_u(k)$, e os coeficientes de projeção de interface
$f_{zm}(k)$, $f_{zp}(k)$ (interpolação linear entre camadas adjacentes),
além dos três coeficientes de extrapolação de contorno de topo
$cf_1,cf_2,cf_3$. Esse conjunto de campos, por depender apenas de
\texttt{config\_ztop}/$n_{VertLevels}$/\texttt{config\_interface\_projection}
-- sem depender da geometria da malha em si -- foi validado como
essencialmente exato (diferenças da ordem de $10^{-6}$ a $10^{-9}$,
atribuíveis a arredondamento de armazenamento em precisão simples) contra
a referência real.

Finalmente, a altura física das interfaces e a métrica vertical são
\begin{equation}
  z_{grid}(k, i) = z_w(k) + A(k)\, h_x(k, i),
  \qquad
  zz(k, i) = \frac{z_w(k{+}1) - z_w(k)}{z_{grid}(k{+}1, i) - z_{grid}(k, i)},
  \label{eq:zgrid_zz}
\end{equation}
e os termos de métrica horizontal $zx_u$ (usados na projeção do gradiente
de altura ao longo das arestas) são calculados por diferença finita
central de $z_{grid}$ entre as duas células vizinhas de cada aresta,
normalizada pela distância entre centros de célula ($dc_{Edge}$).

\subsection{Termos de métrica para advecção de terceira ordem ($zb$, $zb^3$)}

Um segundo achado de implementação, corrigindo uma suposição do plano
original, foi que os campos $zb$/$zb^3$ -- usados no cálculo do vento
vertical diagnóstico da Fase~4 (\cref{cap:fase4}) para o esquema de
advecção de terceira ordem de $\theta$ -- \textbf{não} vêm prontos no
arquivo de malha estática, ao contrário do que se assumia inicialmente;
eles dependem de $z_{grid}$ (produto desta própria fase) e precisam ser
recalculados após ela, usando os coeficientes de segunda derivada
horizontal $\texttt{deriv\_two}$ (esses sim já presentes no arquivo de
malha estática). Uma nova subrotina (\texttt{compute\_zb}, incluída no
mesmo módulo desta fase) foi escrita para esse fim, extraída literalmente
do bloco ``For $z$-metric term in omega equation'' do código de
referência.

\section{Validação numérica}
\label{sec:validacao_fase2}

Após a correção do \cref{achado:hybrid}, a grade vertical calculada foi
comparada campo a campo contra o $z_{grid}$/$zz$ real de um
\texttt{init.nc} de produção (malha \texttt{SouthAmerica}, 17064 células,
56 interfaces). O erro máximo caiu para $\SI{171}{\metre}$ e o erro médio
para $\SI{0,054}{\metre}$. A análise da distribuição espacial do resíduo
confirmou que ele está \textbf{concentrado quase inteiramente} nas sete
camadas do anel de fronteira da malha regional
(\texttt{bdyMaskCell}~$\in\{1,\dots,7\}$, cerca de 2600 das 17064
células): nas células estritamente interiores
(\texttt{bdyMaskCell}$=0$, 82\% da malha), o erro máximo cai para
$\SI{10,9}{\metre}$ -- fisicamente irrelevante. Essa correlação espacial
tem explicação física conhecida: a rotina de referência aplica um passo
de \emph{mistura de terreno de fronteira}
(\texttt{config\_blend\_bdy\_terrain}) \emph{antes} de chamar a geração
da grade vertical, misturando a altura de terreno das células de
fronteira com o terreno do dado de fundo -- passo ainda não implementado
nesta rota (discutido como limitação conhecida no
\cref{cap:discussao}).

\chapter{Fase 3 -- Interpolação vertical}
\label{cap:fase3}

\section{Objetivo e achado sobre a rotina real}

A Fase~3 interpola, para cada coluna vertical da malha de destino, os
campos já remapeados horizontalmente (Fase~1, ainda organizados em
níveis de pressão fixa) para os níveis de modelo da grade vertical
nativa calculada na Fase~2. O plano de implementação original assumia
que essa interpolação reaproveitaria a rotina \texttt{interp\_tofixed\_pressure}
já usada por outro estágio do \selfgrib{} (extração dos campos
isobáricos de diagnóstico). A leitura do código-fonte real do
\texttt{init\_atmosphere\_model} mostrou que essa suposição estava
incorreta: aquela rotina é usada apenas para o \emph{diagnóstico}
isobárico de saída (\texttt{diag.*.nc}), um caminho de código
inteiramente distinto. A interpolação que efetivamente constrói o
estado inicial do modelo usa uma função \textbf{local} de nome
\texttt{vertical\_interp}, definida dentro da própria rotina
\texttt{init\_atm\_case\_gfs} -- uma interpolação linear simples
\textbf{em altura} (não em pressão), ponto a ponto, com extrapolação
configurável nos extremos da coluna.

\section{A função de interpolação}

Dado um perfil de coluna já ordenado por altura crescente,
$\{(z_k, \phi_k)\}_{k=1}^{n}$, e uma altura-alvo $z^\star$
(tipicamente o ponto médio de uma camada de modelo,
$z^\star = \tfrac12[z_{grid}(k) + z_{grid}(k{+}1)]$), o valor
interpolado é
\begin{equation}
  \phi(z^\star) =
  \begin{cases}
    \phi_1 + \left(\dfrac{\phi_2-\phi_1}{z_2-z_1}\right)(z^\star - z_1), & z^\star < z_1 \quad \text{(extrap. linear)} \\[0.8em]
    \phi_1 - 0{,}0065\,(z^\star - z_1), & z^\star < z_1 \quad \text{(extrap. tipo lapso, s\'o T)} \\[0.8em]
    w_-\,\phi_k + w_+\,\phi_{k+1}, \quad w_- = \dfrac{z_{k+1}-z^\star}{z_{k+1}-z_k},\ w_+ = \dfrac{z^\star - z_k}{z_{k+1}-z_k}, & z_k \le z^\star < z_{k+1} \\[0.8em]
    \phi_n + \left(\dfrac{\phi_n-\phi_{n-1}}{z_n-z_{n-1}}\right)(z^\star - z_n), & z^\star \ge z_n
  \end{cases}
  \label{eq:vertical_interp}
\end{equation}
O modo de extrapolação abaixo do primeiro ponto é selecionável por campo:
constante, linear, ou taxa de lapso padrão ($\SI{-0.0065}{\kelvin\per\metre}$,
usada apenas para temperatura, seguindo \texttt{config\_extrap\_airtemp
= 'lapse-rate'} confirmado no \textit{namelist} real de produção). O
parâmetro \texttt{order} presente na assinatura original da função é
recebido mas nunca de fato usado no corpo do código de referência --
a interpolação é sempre linear, independentemente do valor passado --
achado confirmado por leitura direta do código-fonte e preservado
fielmente (isto é, o parâmetro não foi implementado na tradução, já que
não tem efeito na rotina original).

\section{Campos interpolados e particularidades por campo}

\begin{table}[H]
  \centering
  \caption{Campos interpolados verticalmente na Fase~3 e seu modo de
    extrapolação abaixo do topo do dado de origem.}
  \label{tab:campos_fase3}
  \begin{tabular}{@{}lll@{}}
    \toprule
    Campo & Extrapolação & Observação \\
    \midrule
    Temperatura ($T$) & taxa de lapso ($-0{,}0065\,\si{\kelvin\per\metre}$) & \texttt{config\_extrap\_airtemp} \\
    Umidade relativa (RH) & constante & \\
    Umidade específica & constante & valores negativos zerados antes \\
    Altura geopotencial & linear & \\
    Pressão & linear (em $\ln p$) & interpola $\ln p$, depois $\exp(\cdot)$ \\
    Vento normal ($u$, aresta) & constante & perfil já projetado (\cref{eq:vento_projecao}) \\
    Pressão de superfície & linear (em $\ln p$) & alvo $z^\star = z_{grid}(1,i)$ (nível do chão) \\
    \bottomrule
  \end{tabular}
\end{table}

O caso da pressão merece nota: interpola-se $\ln p$, não $p$ diretamente
-- uma escolha do próprio código de referência que garante positividade
do resultado interpolado/extrapolado, já que $\exp(\cdot)$ de qualquer
valor real é sempre positivo, o que não seria garantido interpolando $p$
linearmente e podendo produzir valores negativos em extrapolação. A
pressão de superfície ajustada por diferença de topografia usa a mesma
fonte de dado (perfil de $\ln p$ por célula), mas avaliada no nível mais
baixo da grade de destino em vez de em um ponto médio de camada --
compensando, célula a célula, a diferença de elevação entre a orografia
da malha de origem e o terreno real da malha de destino.

\section{Interpolação do vento}

Conforme descrito na \cref{sec:vento_arestas}, o perfil de vento é
projetado na direção normal da aresta de destino \emph{antes} da
interpolação vertical (\cref{eq:vento_projecao}), usando como coordenada
de altura a média das alturas geopotenciais das duas células vizinhas da
aresta. A \cref{eq:vertical_interp} é então aplicada a esse perfil já
projetado, com extrapolação constante (não linear/lapso, que fariam
menos sentido físico para uma componente de velocidade).

\chapter{Fase 4 -- Balanço hidrostático}
\label{cap:fase4}

\section{Visão geral}

A Fase~4 é a de maior risco numérico de todo o trabalho: transforma os
campos de temperatura, pressão, umidade e vento já interpolados
verticalmente (Fase~3) em um conjunto de variáveis prognósticas
mutuamente consistentes com a métrica vertical e com o balanço
hidrostático de referência do \mpasa{} (\cref{sec:hidrostatico_teoria})
-- densidade seca $\rho_{zz}$, temperatura potencial úmida
$\theta_m$, pressão total $p$, água precipitável, e vento vertical
diagnóstico $w$. O algoritmo, reproduzido aqui passo a passo, foi
extraído literalmente do bloco \texttt{config\_met\_interp} da rotina
\texttt{init\_atm\_case\_gfs} de referência, preservando deliberadamente
duas fórmulas de $\rho_{zz}$ aparentemente redundantes que na verdade não
são -- um ponto de risco identificado explicitamente durante a
implementação e discutido na \cref{sec:duas_formulas}.

\section{Diagnóstico inicial de exner, $\theta$ e $\rho_{zz}$}

A partir da pressão e temperatura já interpoladas (Fase~3), calcula-se,
para cada célula $i$ e nível $k$, a função de Exner
$\Pi = (p/p_0)^{R/c_p}$, converte-se temperatura em temperatura
potencial, e obtém-se uma primeira estimativa \emph{diagnóstica} da
densidade seca:
\begin{equation}
  \Pi(k,i) = \left(\frac{p(k,i)}{p_0}\right)^{R/c_p},
  \qquad
  \theta(k,i) = T(k,i)\left(\frac{p_0}{p(k,i)}\right)^{R/c_p},
  \label{eq:exner_theta}
\end{equation}
\begin{equation}
  \rho_{zz}^{(0)}(k,i) = \frac{1}{1+q_v(k,i)}\cdot
  \frac{p(k,i)}{R\,\Pi(k,i)\,\theta(k,i)\left[1+(R_v/R - 1)\,q_v(k,i)\right]},
  \label{eq:rho_diagnostico}
\end{equation}
com $p_0=\SI{e5}{\pascal}$, $R=\SI{287}{\joule\per\kilo\gram\per\kelvin}$,
$c_p = \tfrac{7}{2}R$, e $R_v/R$ a razão entre as constantes de gás do
vapor d'água e do ar seco. Essa densidade diagnóstica é usada apenas para
a água precipitável (\cref{eq:precipitavel}) e para o acoplamento inicial
com o estado de referência -- \emph{não} é a densidade final do problema,
que exige resolver a relação hidrostática de forma consistente
(\cref{sec:iteracao}).

\begin{equation}
  W = \sum_{k} \rho_{zz}^{(0)}(k,i)\, q_v(k,i)\, \big[z_{grid}(k{+}1,i)-z_{grid}(k,i)\big].
  \label{eq:precipitavel}
\end{equation}

\section{Estado de referência isotérmico seco}

O estado base $(\,\overline{p},\overline{\Pi},\overline{\rho},\overline{\theta}\,)$
é uma atmosfera hipotética isotérmica seca a $T_0=\SI{250}{\kelvin}$,
avaliada no ponto médio geométrico de cada camada $z_m(k,i) =
\tfrac12[z_{grid}(k{+}1,i)+z_{grid}(k,i)]$:
\begin{align}
  \overline{p}(k,i) &= p_0\, \exp\!\left(-\frac{g\, z_m(k,i)}{R\,T_0}\right), \label{eq:estado_base_p}\\
  \overline{\Pi}(k,i) &= \left(\overline{p}(k,i)/p_0\right)^{R/c_p},
  \qquad
  \overline{\rho}(k,i) = \frac{\overline{p}(k,i)}{R\,T_0},
  \qquad
  \overline{\theta}(k,i) = \frac{T_0}{\overline{\Pi}(k,i)}. \label{eq:estado_base_resto}
\end{align}
Em seguida, tanto o estado base $\overline{\rho}$ quanto a densidade
diagnóstica $\rho_{zz}^{(0)}$ são \emph{acoplados à métrica vertical}
(divididos por $zz(k,i)$, \cref{eq:zgrid_zz}), e definem-se as
perturbações de referência
$p'(k,i) = p(k,i) - \overline{p}(k,i)$ e
$\rho'(k,i) = \rho_{zz}^{(0)}(k,i) - \overline{\rho}(k,i)$.

\section{Duas fórmulas de $\rho_{zz}$: por que não são redundantes}
\label{sec:duas_formulas}

\begin{caixaachado}
Um ponto de risco identificado explicitamente durante a implementação --
e preservado deliberadamente, sem tentativa de ``simplificação'' -- é
que o nível mais baixo ($k=1$) e os níveis subsequentes usam
\emph{duas fórmulas distintas} para $\rho_{zz}$, com expoentes e fatores
de correção de umidade diferentes:
\begin{equation}
  \rho_{zz}(1,i) = \frac{1}{zz(1,i)}\cdot\frac{p_0}{R}\cdot
  \frac{\left(p(1,i)/p_0\right)^{c_v/c_p}}{T(1,i)\left[1+1{,}61\,q_v(1,i)\right]},
  \label{eq:rho_nivel1}
\end{equation}
usando o expoente $c_v/c_p$ (não $R/c_p$ como em
\cref{eq:rho_diagnostico}) e o fator de correção de umidade $1{,}61\,q_v$
aplicado diretamente sobre $T$ (não a forma $(R_v/R-1)\,q_v$ do cálculo
diagnóstico inicial). A mesma segunda fórmula -- com o mesmo expoente
$c_v/c_p$ e o mesmo fator $1{,}61\,q_v$ -- reaparece dentro da iteração
hidrostática (\cref{eq:rho_iteracao}) para os níveis $k\ge 2$. Uma
tentativa de ``reconciliar'' essas duas expressões em uma única fórmula
canônica, por parecerem redundantes à primeira vista, introduziria um
erro sistemático sutil no campo de densidade -- por isso ambas foram
preservadas exatamente como aparecem no código-fonte de referência,
documentadas explicitamente no código Fortran resultante para que futuras
manutenções não as \emph{simplifiquem} por engano.
\end{caixaachado}

\section{Solução hidrostática iterativa}
\label{sec:iteracao}

Para os níveis $k=2,\dots,n_{VertLevels}$, resolve-se, célula a célula,
um sistema não linear pela relação hidrostática discreta do \mpasa{},
iterando até convergência (limiar $|\Delta p'| < \SI{0.0001}{\pascal}$,
máximo de 30 iterações):
\begin{align}
  p'(k,i) &\leftarrow p'(k{-}1,i) - \Big[f_{zm}(k)\rho'(k,i) + f_{zp}(k)\rho'(k{-}1,i)\Big]\,g\,\Delta z_u(k) \nonumber\\
  &\qquad - \Big[f_{zm}(k)\rho_{zz}(k,i)q_v(k,i) + f_{zp}(k)\rho_{zz}(k{-}1,i)q_v(k{-}1,i)\Big]\,g\,\Delta z_u(k), \label{eq:pp_iteracao}\\[0.4em]
  p(k,i) &= p'(k,i) + \overline{p}(k,i), \qquad
  \Pi(k,i) = \left(p(k,i)/p_0\right)^{R/c_p}, \label{eq:p_iteracao}\\[0.4em]
  \rho_{zz}(k,i) &= \frac{1}{zz(k,i)}\cdot\frac{p(k,i)}{R\,\Pi(k,i)\,T(k,i)\big[1+1{,}61\,q_v(k,i)\big]},
  \qquad
  \rho'(k,i) = \rho_{zz}(k,i)-\overline{\rho}(k,i). \label{eq:rho_iteracao}
\end{align}
Os coeficientes $f_{zm}(k)$, $f_{zp}(k)$ e $\Delta z_u(k)$ são exatamente
os já calculados na Fase~2 (\cref{cap:fase2}) -- é essa reutilização
direta que garante a consistência entre a métrica vertical e o campo de
massa resultante, mencionada na \cref{sec:hidrostatico_teoria}. Ao final
da convergência em todos os níveis, a temperatura potencial é
transformada de $\theta$ (seca) para $\theta_m$ (potencial úmida,
variável prognóstica real do modelo) por
$\theta_m(k,i) = \theta(k,i)\big[1+1{,}61\,q_v(k,i)\big]$, e a
perturbação $\rho'$ é desacoplada da métrica vertical multiplicando de
volta por $zz(k,i)$.

\section{Vento normal ponderado por densidade e vento vertical diagnóstico}

O produto $ru$ nas arestas -- a variável de fluxo de massa efetivamente
usada pelo núcleo dinâmico, não apenas o vento $u$ -- é
\begin{equation}
  ru(k,e) = u(k,e)\cdot \tfrac12\big[\rho_{zz}(k,c_1(e)) + \rho_{zz}(k,c_2(e))\big].
  \label{eq:ru}
\end{equation}
O vento vertical diagnóstico $w$ é obtido por um balanço de massa: a
divergência horizontal do fluxo $ru$ ao redor de cada célula, projetada
verticalmente pelos termos de métrica $zb$ (\cref{cap:fase2}) e, dado que
o \textit{namelist} real usa \texttt{config\_theta\_adv\_order=3},
acrescida do termo de correção de terceira ordem via $zb^3$ ponderado
pelo coeficiente \texttt{config\_coef\_3rd\_order}~$=0{,}25$:
\begin{equation}
  \rho w\big|_{k} \propto \sum_{e \in \partial i}
  \pm\Big[f_{zm}(k)\,zz(k,i)+f_{zp}(k)\,zz(k{-}1,i)\Big]
  \Big[zb(k,\cdot,e) + \operatorname{sgn}(ru)\cdot 0{,}25\, zb^3(k,\cdot,e)\Big]\, ru_{proj}(k,e),
  \label{eq:rw}
\end{equation}
com o sinal e o índice do segundo argumento de $zb$/$zb^3$ dependendo de
qual das duas células vizinhas da aresta $e$ é a célula $i$ em questão
(convenção idêntica à já usada no cálculo de $zx_u$, Fase~2). O vento
vertical final é $w(k,i) = \rho w|_k \big/ [f_{zp}(k)\rho_{zz}(k{-}1,i) +
f_{zm}(k)\rho_{zz}(k,i)]$, com $w(1,\cdot)=w(n_{VertLevels}+1,\cdot)=0$
por condição de contorno (nível do chão e topo rígido, respectivamente).

\section{Pressão de superfície diagnóstica e campos finais}

Um último campo diagnóstico, $\texttt{surface\_pressure}$ -- distinto do
$\texttt{psfc}$ já calculado na Fase~3 --, é obtido por extrapolação de
segunda ordem a partir dos dois níveis mais baixos:
\begin{equation}
  \texttt{surface\_pressure}(i) = \frac{0{,}5\,g}{rdz_w(1)}
  \Big[1{,}25\,\rho_{zz}(1,i)(1+q_v(1,i)) - 0{,}25\,\rho_{zz}(2,i)(1+q_v(2,i))\Big] + p'(1,i) + \overline{p}(1,i).
  \label{eq:surface_pressure}
\end{equation}
Finalmente, os campos prognósticos de saída são
$\rho(k,i) = \rho_{zz}(k,i)\cdot zz(k,i)$ (densidade física, desacoplada
da métrica) e $\theta(k,i) = \theta_m(k,i)/[1+1{,}61\,q_v(k,i)]$
(restaurando a temperatura potencial seca a partir da úmida, para o
campo \texttt{theta} de saída do \texttt{init.nc}).

\chapter{Fase 6 -- Campos de superfície e solo}
\label{cap:fase6}

\section{Objetivo e achado de organização do código-fonte}

A Fase~6 calcula os campos do \texttt{init.nc} que não fazem parte da
coluna atmosférica propriamente dita -- temperatura profunda do solo,
temperatura de pele corrigida, fração de vegetação e albedo de fundo,
indicadores de neve e gelo marinho, umidade de solo. A categorização
inicial do plano de trabalho assumia que a maioria desses campos seria
uma cópia direta do dado de \textit{first-guess}, sem cálculo adicional.
Essa suposição estava parcialmente incorreta: nenhum desses campos é, de
fato, calculado dentro da mesma rotina que gera o restante do
\texttt{init.nc} (\texttt{init\_atm\_case\_gfs}); eles pertencem a uma
rotina \emph{separada} do código-fonte de referência
(\texttt{physics\_initialize\_real}, em
\texttt{mpas\_atmphys\_initialize\_real.F}), encontrada apenas ao buscar
diretamente pelos nomes das variáveis de saída em todo o diretório
\texttt{core\_init\_atmosphere} -- uma lição de que a estrutura de módulos
do código de referência não segue necessariamente a mesma organização
lógica das fases do problema físico.

\section{Máscara terra/água e correção de temperatura de pele}

A máscara \texttt{xland} é definida diretamente a partir da máscara
estática de uso do solo (\texttt{landmask}, já presente na malha):
$\texttt{xland}=1$ sobre terra, $\texttt{xland}=2$ sobre água. A
temperatura de pele ($\texttt{skintemp}$) é corrigida por um termo de
lapso térmico padrão, compensando a diferença de elevação entre a
orografia usada pela fonte do \textit{first-guess}
($z_{soilhgt}$, campo \texttt{SOILHGT}) e o terreno real da malha de
destino ($z_{ter}$):
\begin{equation}
  T_{skin}^{corrigida}(i) = T_{skin}^{fg}(i) - 0{,}0065\,\big[z_{ter}(i) - z_{soilhgt}(i)\big].
  \label{eq:skintemp}
\end{equation}
A temperatura profunda do solo ($\texttt{tmn}$) usa a mesma taxa de
lapso, mas referenciada à temperatura climatológica anual do solo
(\texttt{soiltemp}, campo estático) sobre terra, e à temperatura de pele
já corrigida sobre água:
\begin{equation}
  T_{mn}(i) =
  \begin{cases}
    T_{soil}^{clim}(i) - 0{,}0065\, z_{ter}(i), & \texttt{landmask}(i)=1 \\
    T_{skin}^{corrigida}(i), & \texttt{landmask}(i)=0
  \end{cases}
  \label{eq:tmn}
\end{equation}

\section{Interpolação climatológica mensal de vegetação e albedo}
\label{sec:monthly_interp}

Os campos estáticos \texttt{greenfrac} e \texttt{albedo12m}, já presentes
no arquivo de malha estática, são climatologias de 12 valores mensais
(um por mês do ano, representando o dia~15 de cada mês). A fração de
vegetação instantânea (\texttt{vegfra}) e o albedo de fundo
(\texttt{sfc\_albbck}) são obtidos por interpolação linear no tempo entre
os dois meses cujo dia~15 está mais próximo da data de início do
experimento, seguindo a mesma convenção do WRF/WPS. Sendo $\mathcal{D}$
a data-alvo (dia juliano do ano) e $\mathcal{D}_1 < \mathcal{D} \le
\mathcal{D}_2$ os dias juliano do dia~15 dos dois meses adjacentes mais
próximos:
\begin{equation}
  \phi(\mathcal{D}) = \frac{\phi_{m_2}\,(\mathcal{D}-\mathcal{D}_1) + \phi_{m_1}\,(\mathcal{D}_2-\mathcal{D})}{\mathcal{D}_2 - \mathcal{D}_1}.
  \label{eq:monthly_interp}
\end{equation}
Os meses de padding usados para cobrir dezembro e janeiro (dias juliano
antes do dia~15 de janeiro ou depois do dia~15 de dezembro) usam um
deslocamento fixo de 31 dias em relação ao dia juliano do mês adjacente
-- não a aritmética exata de calendário do ano anterior/seguinte -- uma
simplificação do próprio código de referência, preservada fielmente.
O albedo de fundo resultante é normalizado por $100$ (a climatologia
estática é armazenada em porcentagem) e forçado a $0{,}08$ sobre células
de água.

\section{Neve, gelo marinho e umidade de solo}

O indicador binário de presença de neve e a profundidade de neve são
\begin{equation}
  \texttt{snowc}(i) =
  \begin{cases} 1, & \texttt{snow}(i) \ge \SI{10}{\kilo\gram\per\meter\squared} \\ 0, & \text{caso contrário} \end{cases},
  \qquad
  \texttt{snowh}(i) = \texttt{snow}(i)\cdot\frac{5{,}0}{1000},
  \label{eq:snow}
\end{equation}
com a razão $5{:}1$ entre lâmina de água equivalente e profundidade de
neve seguindo a mesma convenção padrão do WRF/Noah.
\begin{caixaachado}
Durante a investigação, duas fórmulas distintas para \texttt{snowc}
foram encontradas em pontos diferentes do código-fonte de referência --
uma com limiar $\texttt{snow}>0$, outra com $\texttt{snow}\ge 10$. A
ambiguidade foi resolvida por confronto direto contra o \texttt{init.nc}
real: das 222 células com neve entre $0$ e $\SI{10}{\kilo\gram\per\meter\squared}$
na malha de teste, nenhuma tinha $\texttt{snowc}=1$ no arquivo real,
confirmando o limiar $\ge 10$ como o efetivamente usado em produção.
\end{caixaachado}

O gelo marinho, sob a configuração real de produção
(\texttt{config\_frac\_seaice=.true.}, confirmada apenas após a
generalização do código para ler o \textit{namelist} real -- ver
\cref{sec:generalizacao}), usa um limiar fracionário mais permissivo:
$\texttt{xice}(i) = \texttt{SEAICE}^{fg}(i)$ se $\ge 0{,}02$, senão $0$;
e $\texttt{seaice}(i)=1$ onde $\texttt{xice}(i)>0$. A reclassificação de
células de água muito fria (temperatura de pele abaixo de um limiar
configurável) em gelo/terra, presente no algoritmo de referência
completo, não foi implementada nesta rota -- discutido como limitação
conhecida no \cref{cap:discussao}. A umidade de solo líquida
(\texttt{sh2o}) é $0$ sobre terra e $1$ sobre água (convenção de
saturação total sobre oceano). Os campos de temperatura/umidade das
quatro camadas de solo do esquema Noah (\texttt{tslb}/\texttt{smois})
são, nesta implementação, copiados diretamente do \textit{first-guess}
em vez de reamostrados verticalmente para profundidades-alvo distintas
-- uma simplificação deliberada, justificada e com seu risco documentado
na \cref{cap:discussao}, válida porque tanto a origem quanto o destino
deste \textit{first-guess} usam o mesmo esquema de solo com as mesmas
quatro profundidades.

\chapter{Fase 7 -- Escrita dos arquivos e execução do modelo}
\label{cap:fase7}

\section{Composição do \texttt{init.nc} completo}

O \texttt{init.nc} real do \mpasa{} contém 135 variáveis. Das fases
anteriores (1 a 6), o conjunto de programas deste trabalho calcula e
escreve diretamente 54 delas -- a grade vertical (Fase~2), os campos de
massa e vento (Fases~3--4), e os campos de superfície/solo (Fase~6). As
81 variáveis restantes são de \textbf{cópia direta} do arquivo de malha
estática (\texttt{static.nc}): geometria e conectividade da malha
(coordenadas de célula/aresta/vértice, tabelas de vizinhança), uso do
solo, orografia bruta, e coeficientes numéricos fixos da malha (não
dependentes do dado meteorológico nem do tempo de execução). Essa
composição é feita por mesclagem de arquivos NetCDF
(\texttt{ncks~-A}, utilitário do pacote NCO), não por um terceiro
programa Fortran -- uma escolha de engenharia que evita duplicar, em
Fortran, a leitura e reescrita de 81 campos que já existem prontos, sem
nenhuma transformação necessária.

\begin{figure}[H]
\centering
\begin{tikzpicture}[
    node distance=8mm and 12mm,
    box/.style={draw, rounded corners, minimum height=0.9cm, minimum width=3.2cm, align=center, font=\footnotesize, fill=blue!6},
    arr/.style={-{Stealth[length=2mm]}, thick},
  ]
  \node[box] (static) {\texttt{<REGIÃO>.static.nc}\\(81 campos estáticos)};
  \node[box, right=of static, fill=green!10] (native) {\texttt{gen\_init\_native}\\(Fases 2+3+4+6)};
  \node[box, below=of native, fill=orange!10] (computed) {\texttt{computed\_fields.nc}\\(54 campos calculados)};
  \node[box, right=of computed, fill=red!8] (merge) {\texttt{ncks -A}\\(mesclagem)};
  \node[box, right=of merge] (final) {\texttt{init.nc}\\(135 campos)};

  \draw[arr] (static) -- (native);
  \draw[arr] (native) -- (computed);
  \draw[arr] (static.south) |- ($(merge.west)+(0,0.35)$);
  \draw[arr] (computed) -- (merge);
  \draw[arr] (merge) -- (final);
\end{tikzpicture}
\caption{Composição do \texttt{init.nc} final: os 54 campos calculados
  pelos programas desta rota são mesclados, via \texttt{ncks -A}, com
  os 81 campos de cópia direta do arquivo de malha estática.}
\label{fig:composicao_init}
\end{figure}

O \texttt{lbc.*.nc} tem um esquema mais simples: apenas os sete campos
prognósticos (\texttt{lbc\_qv}, \texttt{lbc\_qc}, \texttt{lbc\_qr},
\texttt{lbc\_u}, \texttt{lbc\_w}, \texttt{lbc\_rho}, \texttt{lbc\_theta}),
um arquivo por tempo de fronteira, sem geometria nenhuma -- a malha e a
grade vertical são lidas pelo \texttt{mpas\_atmosphere} do próprio
\texttt{init.nc}, via o \emph{stream} de entrada. O programa
\texttt{gen\_lbc\_native} reaproveita, por isso, a grade vertical já
calculada e gravada no \texttt{init.nc} (não a recalcula), rodando apenas
as Fases~3--4 para cada tempo requisitado.

\section{Dois bugs de formato encontrados apenas na execução real}
\label{sec:bugs_formato}

A validação numérica campo a campo contra arquivos de referência
(\cref{cap:resultados}), por mais rigorosa que tenha sido, não é
suficiente para garantir que um arquivo NetCDF seja de fato
\emph{consumível} pelo programa de destino -- um ponto que se revelou
concretamente ao tentar, pela primeira vez, alimentar o
\texttt{mpas\_atmosphere} real com os arquivos gerados por esta rota.

\begin{caixaachado}
\textbf{Ausência da variável \texttt{xtime}.} A primeira tentativa de
execução falhou imediatamente, antes mesmo do primeiro passo de tempo,
com o erro \texttt{"ERROR: File SouthAmerica.init.nc does not contain
the time 2026-01-01\_00:00:00"}. O \texttt{init.nc} gerado continha a
variável \texttt{initial\_time} (uma string informativa, escrita
corretamente), mas não a variável \texttt{xtime(Time, StrLen)} -- que é,
na verdade, a que o \texttt{mpas\_atmosphere} lê para confirmar que o
arquivo de entrada contém o tempo solicitado pelo
\texttt{config\_start\_time}. A variável havia sido listada, na fase de
levantamento inicial das 135 variáveis do \texttt{init.nc} real, mas não
percebida como funcionalmente distinta de \texttt{initial\_time} até
esse ponto -- um lembrete de que um levantamento de nomes de variável não
substitui saber \emph{qual delas o consumidor realmente lê}.

\textbf{Lixo de memória no \texttt{xtime} do \texttt{lbc.*.nc}.} Corrigido
o problema acima, uma segunda falha ocorreu já dentro do primeiro passo
de tempo, com um erro de execução do próprio \emph{runtime} Fortran
(\texttt{"Fortran runtime error: Bad integer for item 1 in list input"})
dentro da rotina de gerenciamento de tempo do framework
(\texttt{mpas\_timekeeping.F}), derrubando simultaneamente todos os 32
processos MPI. A causa, encontrada inspecionando os \emph{bytes} brutos
do arquivo gerado, foi uma chamada \texttt{nf90\_put\_var} recebendo uma
expressão de string de 24 caracteres para um campo NetCDF declarado com
64 bytes: os 40 bytes restantes ficavam preenchidos com lixo de memória
não inicializado (confirmado como binário arbitrário, não espaços em
branco), que o \emph{parser} de data do \textit{framework} -- uma leitura
formatada por posição fixa de coluna, rígida por natureza -- não
conseguia interpretar. A correção foi declarar explicitamente uma
variável Fortran do tipo \texttt{character(len=64)} antes de passá-la ao
\texttt{nf90\_put\_var} (o Fortran preenche automaticamente o excedente
com espaços em branco ao atribuir uma string mais curta a uma variável de
tamanho fixo maior), em vez de passar a expressão concatenada
diretamente.
\end{caixaachado}

Um terceiro problema, de natureza puramente operacional (não um bug do
código Fortran), também surgiu apenas na execução real: o script de
orquestração da rodada carrega um arquivo de ambiente do cluster que
executa \texttt{conda deactivate} internamente; quando o script é
invocado como \texttt{bash script.bash} (\emph{shell} não interativo), o
\texttt{\textasciitilde/.bashrc} não é lido por padrão mesmo com o
\textit{hook} do \texttt{conda} já inicializado uma vez no terminal
interativo do usuário, e o comando falha, derrubando o script inteiro
(\texttt{set -e}). A correção carregou defensivamente o \textit{hook} do
\texttt{conda} (\texttt{etc/profile.d/conda.sh}) no início do script, se
necessário, antes de delegar ao arquivo de ambiente do cluster.

A lição transversal aos três problemas: nenhum deles era de física ou de
numérica -- o \texttt{init.nc}/\texttt{lbc.*.nc} já batiam campo a campo
com a referência antes de qualquer um deles ser encontrado -- e mesmo
assim o modelo real recusava ou travava com os arquivos. Validação
numérica contra referência não substitui validar contra o consumidor
real; um arquivo pode estar fisicamente correto e ainda assim ser
inutilizável por um detalhe de formato invisível a uma comparação
campo a campo.

\section{Execução completa}

Corrigidos os três problemas, a execução do \texttt{mpas\_atmosphere}
real, com 32 processos MPI, a partir do \texttt{init.nc}/\texttt{lbc.*.nc}
gerados inteiramente por esta rota, completou 240 passos de tempo (24
horas de integração, malha \texttt{SouthAmerica}, $\sim\SI{60}{\kilo\metre}$
de resolução) em aproximadamente 106 segundos de tempo de parede, sem
nenhuma mensagem de erro ou erro crítico (apenas 24 avisos inofensivos,
relativos a variáveis de gelo/graupel/neve ausentes do
\textit{first-guess} -- esperado, já que o dado de origem não modela
esses hidrometeoros). A condição de contorno lateral foi aplicada
corretamente nos quatro intervalos de seis horas cobertos pelos arquivos
\texttt{lbc.*.nc}, gerando 25 arquivos de saída
(\texttt{history.*.nc}, um por hora). Os resultados quantitativos dessa
execução, e sua comparação com a rodada de referência gerada pela rota
original, são apresentados no \cref{cap:resultados}.

\chapter{Implementação em Fortran}
\label{cap:implementacao}

\section{Princípio norteador: extração literal, não reformulação}

A convenção adotada em todo o desenvolvimento -- já estabelecida em
partes anteriores do projeto \selfgrib{} e mantida aqui -- foi a de
\textbf{extração literal} das fórmulas e algoritmos do código-fonte de
referência do MPAS-Model, em vez de reimplementá-los ``do zero'' a partir
da descrição em artigos ou manuais. A justificativa é dupla: primeiro,
garante consistência numérica exata com o que o modelo real produz,
eliminando uma classe inteira de divergência que surgiria de decisões de
implementação sutilmente diferentes (ordem de operações, escolha entre
formas matematicamente equivalentes de uma mesma expressão, tratamento de
casos de borda); segundo, torna a origem de cada linha de código
rastreável e auditável -- cada módulo novo carrega, em comentário de
cabeçalho, a referência exata do arquivo/rotina/faixa de linhas do
código-fonte original de onde foi extraído, e a versão específica
consultada (\texttt{mpas-bundle-3.0.2}, confirmada como o fonte exato que
gerou os binários de produção usados neste projeto, distinta e por vezes
divergente da \emph{branch} principal do repositório público -- por
exemplo, a fórmula do fator de suavização $s(i)$ na
\cref{eq:suavizacao_hx} usa $z_t$ como referência de altura na versão
3.0.2, mas $z_h$ na \emph{branch} \texttt{master} mais recente).

\section{Estrutura de módulos}

A \cref{tab:modulos} resume os módulos Fortran novos, organizados pela
fase a que pertencem, com suas dependências e origem no código-fonte de
referência.

\begin{table}[H]
\centering
\small
\caption{Módulos Fortran novos desenvolvidos neste trabalho.}
\label{tab:modulos}
\begin{tabular}{@{}p{3.1cm}p{1.3cm}p{6.7cm}@{}}
\toprule
Módulo/programa & Fase & Origem no código-fonte de referência \\
\midrule
\texttt{hinterp\_native.F90} & 1 & Cópia estrutural de \texttt{convert\_mpas.F90}, reutilizando sem modificação \texttt{mpas\_mesh}, \texttt{target\_mesh}, \texttt{remapper} \\
\texttt{vertical\_grid.F90} & 2 & \texttt{init\_atm\_case\_gfs}, bloco \texttt{config\_tc\_vertical\_grid} \\
\texttt{vinterp\_native.F90} & 3 & Função local \texttt{vertical\_interp} de \texttt{init\_atm\_case\_gfs} \\
\texttt{hydrostatic.F90} & 4 & \texttt{init\_atm\_case\_gfs}, bloco \texttt{config\_met\_interp} \\
\texttt{surface\_fields.F90} & 6 & \texttt{mpas\_atmphys\_initialize\_real.F}, \texttt{mpas\_atmphys\_date\_time.F} \\
\texttt{gen\_init\_native.F90} & 2+3+4+6+7 & Programa-driver, novo (orquestração + escrita NetCDF) \\
\texttt{gen\_lbc\_native.F90} & 3+4+7 & Programa-driver, novo (reusa grade vertical do \texttt{init.nc}) \\
\texttt{namelist\_config.F90} & todas & Novo (leitura de \textit{namelist} real, sem equivalente direto) \\
\bottomrule
\end{tabular}
\end{table}

\section{Convenção de precisão numérica: uma exceção deliberada}

Todo o restante do projeto \selfgrib{} usa precisão dupla
(\texttt{RKIND}, \texttt{real*8}) como padrão. A única exceção
deliberada é o bloco de suavização iterativa de terreno da Fase~2
(\cref{sec:suavizacao}): o log real de produção revela que o binário do
\texttt{init\_atmosphere\_model} foi compilado com \texttt{"Default real
precision: single"} -- isto é, precisão simples (\texttt{real*4}) em toda
a base de código, não apenas nesse bloco. Como o critério de aceitação
por limiar desse algoritmo (\cref{eq:suavizacao_hx}) é sensível a
diferenças mínimas de arredondamento, compostas ao longo de até 85
passadas por até 53 níveis, as variáveis diretamente envolvidas nessa
suavização (terreno, $h_x$, acumuladores) foram declaradas com um
parâmetro de precisão simples dedicado
(\texttt{SPKIND=selected\_real\_kind(6)}), propositalmente, para casar
com o binário real -- uma escolha investigada e confirmada
experimentalmente (\cref{achado:hybrid}) como \emph{não} sendo, sozinha,
a causa da divergência numérica original, mas mantida por rigor mesmo
após a causa real ter sido encontrada, já que reduz ruído residual de
comparação.

\section{Leitura do \textit{namelist} real em tempo de execução}
\label{sec:namelist_config}

O módulo \texttt{namelist\_config.F90} define um tipo derivado
(\texttt{init\_atm\_config\_type}) que espelha os campos de configuração
efetivamente usados pelas fases deste trabalho, com valores-padrão
sensatos, e uma rotina (\texttt{read\_init\_atm\_namelist}) que lê, via
bloco \texttt{NAMELIST} nativo do Fortran, os grupos
\texttt{\&nhyd\_model}, \texttt{\&dimensions}, \texttt{\&data\_sources},
\texttt{\&vertical\_grid}, \texttt{\&interpolation\_control} e
\texttt{\&preproc\_stages} do arquivo \texttt{namelist.init\_atmosphere}
real do experimento -- o mesmo arquivo já consumido pela rota original,
sem nenhum formato novo introduzido. A leitura é tolerante a grupos ou
variáveis ausentes de um \textit{namelist} específico (por exemplo,
\texttt{config\_fg\_interval} só existe no \textit{namelist} do
\texttt{lbc\_run}, não no do \texttt{init\_run}): os valores-padrão do
tipo derivado são inicializados antes da leitura, de modo que uma
variável ausente simplesmente mantém seu padrão em vez de ficar
indefinida.

\begin{lstlisting}[caption={Trecho ilustrativo de \texttt{namelist\_config.F90}: leitura tolerante de múltiplos grupos de \textit{namelist}.}, label={lst:namelist}]
type :: init_atm_config_type
    real (kind=RKIND) :: config_ztop = 30000.0_RKIND
    integer            :: config_nsm  = 30
    real (kind=RKIND) :: config_dzmin = 0.3_RKIND
    logical            :: config_blend_bdy_terrain = .false.
    logical            :: config_frac_seaice       = .true.
    character (len=64) :: config_start_time = ''
    ! ... demais campos config_*
end type init_atm_config_type

subroutine read_init_atm_namelist(filename, cfg)
    ! ...
    open(newunit=unit_nml, file=trim(filename), status='old', action='read')
    read(unit_nml, nml=nhyd_model,    iostat=ios); rewind(unit_nml)
    read(unit_nml, nml=vertical_grid, iostat=ios); rewind(unit_nml)
    ! ... demais grupos, cada um tolerante a ausencia
    close(unit_nml)
    call parse_config_start_time(cfg)
end subroutine
\end{lstlisting}

Essa generalização, aplicada de forma consistente aos dois programas
principais (\texttt{gen\_init\_native}, \texttt{gen\_lbc\_native}),
resultou em uma constatação importante de validação: recompilar e
re-executar toda a cadeia com os parâmetros agora lidos dinamicamente,
em vez de fixos, produziu um resultado \textbf{numericamente idêntico,
bit a bit}, ao obtido com os valores anteriormente fixos no
código-fonte -- confirmando que a generalização não alterou nenhum
cálculo, apenas a origem dos parâmetros.

\section{Reúso de dependências e organização do build}

O sistema de \textit{build} (\texttt{Makefile}) mantém a convenção já
estabelecida no projeto \selfgrib{} de compilar \texttt{convert\_mpas}
como dependência externa antes do restante da pipeline, já que
\texttt{hinterp\_native} (Fase~1) reutiliza diretamente seus módulos
objeto já compilados (\texttt{mpas\_mesh.o}, \texttt{target\_mesh.o},
\texttt{remapper.o}, \texttt{scan\_input.o}), sem recompilá-los. Os
demais módulos novos (Fases~2 a~7) dependem apenas de
\texttt{mpas\_kind\_types.F90} e \texttt{mpas\_constants.F90}, ambos já
existentes no projeto -- nenhuma dependência externa nova foi introduzida
por este trabalho, além do utilitário NCO (\texttt{ncks}) usado na
composição final do \texttt{init.nc} (\cref{cap:fase7}), tipicamente já
disponível em ambientes de HPC voltados a geociências.

\chapter{Resultados}
\label{cap:resultados}

Este capítulo consolida os resultados quantitativos de validação
apresentados ao longo dos capítulos anteriores, organizados por nível de
integração: (i)~validação isolada da Fase~1 (propriedade de exatidão);
(ii)~validação campo a campo do \texttt{init.nc} completo contra um
arquivo real de produção; (iii)~validação do \texttt{lbc.*.nc}, por
auto-consistência e por comparação com a referência real; e
(iv)~validação funcional, comparando a previsão de 24 horas obtida a
partir dos arquivos desta rota com a previsão de referência gerada pela
rota original. Todos os números desta seção referem-se ao caso de estudo
usado consistentemente em todo o desenvolvimento: malha regional
\texttt{SouthAmerica} (recorte de $x1.163842$, 17064 células, 51648
arestas, $\sim\SI{60}{\kilo\metre}$ de resolução nominal), rodada iniciada
em \texttt{2026-01-01\_00:00:00}.

\section{Validação da interpolação horizontal (Fase 1)}

Conforme discutido na \cref{sec:validacao_fase1}, a interpolação de uma
malha de teste (2562 células) para os próprios centros de célula dessa
mesma malha (\emph{auto-\emph{scatter}}) produziu erro identicamente nulo
em 100\% das células, tanto para os valores de campo quanto para as
coordenadas de saída -- confirmando a implementação correta da
propriedade de exatidão nos vértices da interpolação baricêntrica
(\cref{sec:baricentrica}).

\section{Validação do \texttt{init.nc} completo}

\begin{table}[H]
\centering
\caption{Erro absoluto campo a campo entre o \texttt{init.nc} gerado por
  esta rota e o \texttt{init.nc} real de produção (17064 células, malha
  \texttt{SouthAmerica}, mesmo tempo inicial). Compilado e executado no
  build/ambiente real de produção (cluster Jaci), não apenas em ambiente
  de desenvolvimento local.}
\label{tab:validacao_init}
\begin{tabular}{@{}lrrr@{}}
\toprule
Campo & Faixa real (mín/máx) & Erro médio & Erro máximo \\
\midrule
$\rho$ [\si{\kilo\gram\per\meter\cubed}] & $0{,}018$ / $1{,}265$ & $1{,}5\times10^{-4}$ & $1{,}84\times10^{-2}$ \\
$\theta$ [\si{\kelvin}] & $273{,}1$ / $825{,}0$ & $0{,}114$ & $3{,}56$ \\
$q_v$ [\si{\kilo\gram\per\kilo\gram}] & $0$ / $0{,}0192$ & $5{,}1\times10^{-5}$ & $2{,}52\times10^{-3}$ \\
$u$ [\si{\meter\per\second}] & $-76$ / $77$ & $8{,}5\times10^{-2}$ & $4{,}47$ \\
$w$ [\si{\meter\per\second}] & $-1{,}52$ / $1{,}08$ & $4{,}0\times10^{-3}$ & $1{,}63$ \\
\texttt{surface\_pressure} [\si{\pascal}] & $5{,}76\times10^4$ / $1{,}03\times10^5$ & $5{,}9$ & $1792$ \\
\texttt{precipw} [\si{\milli\meter}] & $0{,}92$ / $62{,}1$ & $0{,}31$ & $4{,}46$ \\
\texttt{skintemp} [\si{\kelvin}] & $273{,}0$ / $307{,}5$ & $0{,}21$ & $7{,}07$ \\
\texttt{tmn} [\si{\kelvin}] & $266{,}2$ / $302{,}1$ & $0{,}15$ & $6{,}95$ \\
\texttt{vegfra} [\si{\percent}] & $0$ / $90{,}6$ & $7{,}3\times10^{-7}$ & $1{,}1\times10^{-5}$ \\
\bottomrule
\end{tabular}
\end{table}

Campos exatos ou essencialmente exatos (erro nulo ou de arredondamento de
armazenamento, $<10^{-6}$): \texttt{xland}, \texttt{xice},
\texttt{seaice}, \texttt{sh2o}, \texttt{dz}, \texttt{u\_init},
\texttt{v\_init}, \texttt{qv\_init}, \texttt{t\_init}, \texttt{qc},
\texttt{qr}, \texttt{sfc\_albbck}, \texttt{dzs}, \texttt{zs}. O arquivo
final contém a totalidade das 135 variáveis do \texttt{init.nc} real
(mais uma variável diagnóstica auxiliar, \texttt{psfc}, sem
correspondente no arquivo real e sem uso pelo núcleo do modelo). Os
resultados na \cref{tab:validacao_init} foram obtidos compilando e
executando o pipeline diretamente no ambiente de produção (compilador
GNU~13.3.0, bibliotecas NetCDF/PIO do cluster), reproduzindo,
\textbf{dígito a dígito}, os mesmos valores já obtidos em ambiente de
desenvolvimento local -- uma confirmação adicional de portabilidade e
robustez numérica da implementação frente a diferenças de compilador e
biblioteca.

\section{Validação do \texttt{lbc.*.nc}}

A auto-consistência entre o \texttt{init.nc} e um \texttt{lbc.*.nc}
gerado para o \emph{mesmo} tempo inicial (mesma malha vertical, mesmo
\textit{first-guess}) foi verificada como \textbf{exata}: diferença
identicamente nula em $\rho$, $\theta$, $q_v$, $u$ e $w$ -- o resultado
esperado, já que ambos os programas executam exatamente o mesmo cálculo
de Fases~3--4 para esse tempo.

\begin{table}[H]
\centering
\caption{Erro absoluto médio/máximo entre o \texttt{lbc.*.nc} gerado por
  esta rota e os arquivos reais de produção, nos cinco tempos de
  fronteira do experimento (intervalo de 6~h).}
\label{tab:validacao_lbc}
\begin{tabular}{@{}lrrrr@{}}
\toprule
Tempo & $u$ [\si{\meter\per\second}] & $\rho$ & $\theta$ [\si{\kelvin}] & $w$ [\si{\meter\per\second}] \\
\midrule
$t_0$ (00h)  & $0{,}086$ / $4{,}47$  & $0{,}0056$ / $0{,}043$ & $0{,}114$ / $3{,}56$ & $0{,}004$ / $1{,}63$ \\
$t_0+6$h  & $0{,}062$ / $5{,}69$  & $0{,}0056$ / $0{,}044$ & $0{,}055$ / $2{,}86$ & $0{,}004$ / $1{,}59$ \\
$t_0+12$h & $0{,}055$ / $7{,}10$  & $0{,}0056$ / $0{,}041$ & $0{,}051$ / $3{,}15$ & $0{,}004$ / $1{,}60$ \\
$t_0+18$h & $0{,}052$ / $12{,}2$  & $0{,}0057$ / $0{,}042$ & $0{,}060$ / $4{,}65$ & $0{,}004$ / $1{,}45$ \\
$t_0+24$h & $0{,}055$ / $12{,}5$  & $0{,}0057$ / $0{,}042$ & $0{,}059$ / $2{,}91$ & $0{,}004$ / $1{,}46$ \\
\bottomrule
\end{tabular}
\end{table}

O erro médio permanece pequeno e estável ao longo dos cinco tempos; o
erro máximo pontual cresce moderadamente nos tempos mais tardios --
compatível com divergência acumulada de previsão real ao longo de até 24
horas somada ao resíduo de fronteira já conhecido (\cref{cap:discussao}),
não com instabilidade numérica. O único problema identificado nesta
comparação não é desta implementação: o campo \texttt{lbc\_qv} do
arquivo de referência real é \textbf{constante em todos os 55 níveis
verticais por célula} -- fisicamente implausível para um campo de
umidade específica, e inconsistente com o comportamento correto (variável
por nível) dos demais campos do mesmo arquivo --, sugerindo um problema
já existente na cadeia de geração de referência (rota via WPS), não
atribuível a esta rota.

\section{Validação funcional: execução do modelo e comparação de previsão}

A previsão de 24 horas gerada a partir dos arquivos desta rota
(\cref{cap:fase7}) foi comparada, campo a campo e por tempo, com a
previsão de referência (mesma malha, mesmo período, gerada pela rota
original via WPS).

\begin{table}[H]
\centering
\caption{Erro absoluto médio/máximo entre a previsão gerada a partir dos
  arquivos desta rota e a previsão de referência, em três tempos de
  saída ao longo da integração de 24~h.}
\label{tab:validacao_previsao}
\begin{tabular}{@{}lrrrrr@{}}
\toprule
Tempo de previsão & $\theta$ [\si{\kelvin}] & $\rho$ & $q_v$ [\si{\kilo\gram\per\kilo\gram}] & $u$ [\si{\meter\per\second}] & $w$ [\si{\meter\per\second}] \\
\midrule
$+0$h  & $0{,}11$ / $3{,}6$ & $1{,}5\times10^{-4}$ / $1{,}8\times10^{-2}$ & $5{,}1\times10^{-5}$ / $2{,}5\times10^{-3}$ & $0{,}08$ / $4{,}5$ & $0{,}004$ / $1{,}6$ \\
$+12$h & $0{,}27$ / $6{,}4$ & $9{,}2\times10^{-4}$ / $2{,}9\times10^{-2}$ & $4{,}7\times10^{-4}$ / $1{,}8\times10^{-2}$ & $0{,}44$ / $9{,}2$ & $0{,}007$ / $0{,}33$ \\
$+24$h & $0{,}41$ / $9{,}8$ & $1{,}1\times10^{-3}$ / $3{,}0\times10^{-2}$ & $6{,}4\times10^{-4}$ / $1{,}8\times10^{-2}$ & $0{,}58$ / $12{,}0$ & $0{,}010$ / $0{,}65$ \\
\bottomrule
\end{tabular}
\end{table}

O padrão observado -- erro pequeno na inicialização, crescendo de forma
suave e limitada ao longo da integração, sem nenhum salto abrupto ou
sinal de instabilidade (ausência total de valores \texttt{NaN} em
qualquer campo, em qualquer tempo) -- é o esperado para duas trajetórias
vizinhas de um sistema dinâmico não linear sensível a condições iniciais
(divergência tipo Lyapunov): as duas rotas partem de condições iniciais
ligeiramente diferentes (pelo resíduo residual documentado nos
Capítulos~\ref{cap:fase2} e~\ref{cap:fase6}) e evoluem, sob a mesma
física e a mesma dinâmica, para estados progressivamente mais distantes
-- não uma divergência causada por um erro estrutural na implementação,
que se manifestaria como blowup numérico, quebra de conservação, ou
surgimento de valores não físicos, nenhum dos quais foi observado. O
próprio log de execução do modelo confirma a sanidade física ao longo de
toda a integração: os valores globais de mínimo/máximo de $w$ e $u$,
impressos a cada passo de tempo, permanecem dentro de faixas
meteorologicamente plausíveis (por exemplo, $w \in [-0{,}61,
+2{,}22]\,\si{\meter\per\second}$, $u \in [\pm 65, \pm 77]\,\si{\meter\per\second}$)
durante toda a janela de 24 horas.

\chapter{Discussão e limitações conhecidas}
\label{cap:discussao}

Este capítulo consolida, de forma explícita, as simplificações
deliberadas e limitações conhecidas da implementação atual -- distintas
de bugs (já corrigidos ao longo dos capítulos anteriores), no sentido de
que são escolhas de escopo documentadas, com risco avaliado e aceito
para o caso de estudo validado, mas que podem se tornar relevantes em
outros contextos de uso.

\section{Mistura de terreno de fronteira não implementada}
\label{sec:blend_bdy_terrain}

A limitação de maior impacto numérico mensurável identificada é a
ausência do passo \texttt{config\_blend\_bdy\_terrain}: no algoritmo de
referência, antes de gerar a grade vertical (\cref{cap:fase2}), a altura
de terreno das células do anel de fronteira da malha regional
(\texttt{bdyMaskCell}~$\ge 1$) é misturada com o terreno do dado de
\textit{fundo} (\textit{first-guess}), suavizando a transição entre o
domínio regional e a condição de contorno que o alimentará. Sem esse
passo, o resíduo de $\sim\SI{171}{\metre}$ documentado na
\cref{sec:validacao_fase2} permanece concentrado exatamente nessa faixa
de células -- que é, por definição, onde a fronteira lateral atua. Para
o caso de validação usado neste relatório, esse resíduo não impediu uma
integração de 24 horas fisicamente sã (\cref{cap:resultados}), mas é a
lacuna mais provável de precisar de atenção em usos futuros que dependam
de maior precisão perto da borda do domínio regional, ou em domínios com
relevo mais acentuado próximo à fronteira do que o caso testado.

\section{Reamostragem vertical de solo simplificada}

O algoritmo de referência reamostra os perfis de temperatura e umidade
de solo do \textit{first-guess} para as quatro profundidades-padrão do
esquema Noah usando uma interpolação linear por profundidade própria
(distinta e independente da grade vertical atmosférica). Esta
implementação copia esses campos diretamente do \textit{first-guess}
sem essa reamostragem -- uma simplificação segura apenas porque, no caso
de uso deste projeto, tanto a origem quanto o destino do
\textit{first-guess} são o próprio \mpasa{}, com o mesmo esquema Noah e
as mesmas quatro profundidades de solo. Essa simplificação deixa de ser
válida se o \textit{first-guess} vier de uma fonte com profundidades de
solo diferentes (por exemplo, GFS via WPS/\texttt{ungrib}).

\section{Reclassificação de gelo marinho não implementada}

O algoritmo de referência reclassifica células de água muito fria
(temperatura de pele abaixo de um limiar configurável,
\texttt{config\_tsk\_seaice\_threshold}) como gelo marinho ou terra,
ajustando em cascata o tipo de uso do solo, textura de solo, albedo
máximo de neve, e os perfis de temperatura/umidade de solo dessas
células. Essa reclassificação não foi implementada. A omissão é
irrelevante para o caso de validação (América do Sul: o campo
\texttt{SEAICE} do \textit{first-guess} é identicamente zero em toda a
malha, confirmado diretamente contra o dado real), mas se tornaria
relevante em qualquer malha regional que inclua regiões de alta
latitude com gelo marinho real.

\section{Umidade específica via umidade relativa não implementada}

O código de referência oferece dois caminhos para diagnosticar a razão
de mistura de vapor d'água: a partir da umidade específica do
\textit{first-guess} (\texttt{config\_use\_spechumd=.true.}) ou, na
ausência desse campo, a partir da umidade relativa via a função de
saturação de Flatau/Thompson (\texttt{rslf}). Apenas o primeiro caminho
foi implementado -- suficiente porque todos os casos de estudo deste
projeto, até o momento, usam \texttt{config\_use\_spechumd=.true.} com um
\textit{first-guess} que de fato contém esse campo. O programa detecta
explicitamente a configuração oposta e encerra com uma mensagem de erro
clara, em vez de produzir um resultado silenciosamente incorreto.

\section{Escopo de validação}

Todo o processo de validação numérica e funcional apresentado neste
relatório usa uma única malha regional (\texttt{SouthAmerica},
$\sim\SI{60}{\kilo\metre}$), um único caso de estudo (24 horas de
integração a partir de \texttt{2026-01-01\_00:00:00}), e
$n_{VertLevels}=55$. A lista de 61 níveis de pressão fixa usada como
representação intermediária entre a extração de campos e a interpolação
horizontal nativa (\texttt{pressure\_levels.F90}, herdada sem
modificação da rota original) é derivada de medianas estatísticas de
altura de nível de modelo -- não específica da América do Sul, e portanto
esperada como razoável em outras malhas -- mas nunca de fato testada em
uma malha/resolução distinta. A generalização estrutural do código
(\cref{sec:generalizacao}) remove a maior barreira à reprodução em outros
casos, mas não substitui uma validação explícita nesse cenário mais
amplo.

\section{Um interpolante linear por partes: limite de suavidade}

A interpolação baricêntrica adotada (\cref{sec:baricentrica}) é
$C^0$ -- contínua entre triângulos adjacentes do dual de Delaunay, mas
com derivada descontínua nas arestas desses triângulos. Para o problema
resolvido aqui (gerar uma condição inicial/de fronteira, consumida
integralmente pelo núcleo dinâmico do modelo, que já aplica sua própria
suavização/filtro numérico durante a integração), essa limitação não se
manifestou como problema na validação funcional (\cref{cap:resultados}).
Ela poderia se tornar relevante caso o mesmo motor de interpolação fosse
reaproveitado para outro uso que dependa diretamente de gradientes
horizontais suaves do campo interpolado -- caso em que os interpolantes
de maior continuidade discutidos na revisão de literatura
(\cref{cap:literatura}, Hiyoshi e Sugihara~\cite{hiyoshi2000continuity})
seriam um candidato natural de extensão.

\section{Síntese das limitações}

\begin{table}[H]
\centering
\caption{Síntese das simplificações/limitações conhecidas, seu impacto
  no caso validado, e a condição sob a qual se tornariam relevantes.}
\label{tab:limitacoes}
\small
\begin{tabular}{@{}p{4.3cm}p{3.6cm}p{5.2cm}@{}}
\toprule
Limitação & Impacto no caso validado & Torna-se relevante quando \\
\midrule
Sem mistura de terreno de fronteira & Resíduo de $\sim\SI{171}{m}$ concentrado no anel de fronteira & Precisão perto da borda regional importa; relevo acentuado na fronteira \\
Solo copiado sem reamostragem vertical & Nenhum (mesma malha/esquema na origem e destino) & \textit{First-guess} com esquema de solo/profundidades distintas \\
Sem reclassificação de gelo marinho & Nenhum (sem gelo marinho no caso) & Malha regional em alta latitude \\
Sem \texttt{qv} via RH & Nenhum (\texttt{config\_use\_spechumd} sempre \texttt{.true.} nos casos testados) & \textit{First-guess} sem umidade específica direta \\
Validação em 1 malha/1 caso & -- & Uso em outra resolução/domínio/estação do ano \\
Interpolação $C^0$ (baricêntrica linear) & Nenhum observado & Uso que dependa de gradiente horizontal suave do campo interpolado \\
\bottomrule
\end{tabular}
\end{table}

\chapter{Conclusão}
\label{cap:conclusao}

\section{Síntese}

Este trabalho demonstrou que é possível eliminar inteiramente a etapa de
reamostragem em grade latitude-longitude regular do processo de geração
de condição inicial e de fronteira lateral do \mpasa{}, interpolando
diretamente entre a malha de Voronoi nativa de origem (uma rodada global
já concluída) e a malha de Voronoi nativa de destino (o domínio
regional), e escrevendo os arquivos \texttt{init.nc}/\texttt{lbc.*.nc}
consumidos pelo núcleo dinâmico \texttt{mpas\_atmosphere} sem depender do
\texttt{init\_atmosphere\_model} original. A viabilidade dessa rota não
era óbvia \emph{a priori}: a investigação do código-fonte de referência
mostrou que o formato binário intermediário do WPS é, por definição, uma
grade regular, de modo que nenhum ``escopo mínimo'' preservaria o ganho
de precisão da interpolação nativa sem substituir inteiramente o
caminho -- horizontal, vertical, hidrostático e de superfície -- por
implementação própria.

O método central -- interpolação baricêntrica na malha dual de Delaunay
-- não é uma proposta original deste trabalho, mas a aplicação de uma
técnica já validada pela comunidade (via o sistema de assimilação
MPAS-DART) a um novo contexto de uso. A contribuição prática está na
extração cuidadosa, tradução literal e validação exaustiva de cada
fórmula fisicamente necessária do código-fonte de referência do
MPAS-Model -- um processo de engenharia reversa disciplinado que revelou,
ao longo do caminho, tanto suposições incorretas do planejamento inicial
quanto bugs reais de implementação, todos documentados com sua
investigação e correção nos capítulos técnicos deste relatório.

Um padrão metodológico se repete em praticamente todos os problemas mais
difíceis encontrados: a causa raiz raramente estava onde o sintoma
aparecia. A divergência de $\sim\SI{2}{\kilo\metre}$ na grade vertical
não estava no algoritmo de suavização (onde o sintoma se manifestava),
mas em uma suposição incorreta sobre um valor padrão de configuração
calculado dezenas de linhas antes. Os dois bugs que impediam a execução
real do modelo não eram de física nem de numérica -- eram de formato de
uma única variável de metadado (\texttt{xtime}) inicialmente considerada
secundária. Em ambos os casos, uma validação aparentemente completa
(revisão de código linha a linha; comparação numérica campo a campo)
não foi suficiente até que a validação fosse estendida às dependências
reais do problema -- o código \emph{upstream} do trecho suspeito, no
primeiro caso; o consumidor final de verdade, no segundo.

\section{Estado atual e alcance da validação}

Ao final deste trabalho, a rota nativa \emph{voronoi-to-voronoi} está
validada em duas camadas independentes: numericamente, campo a campo,
contra arquivos \texttt{init.nc}/\texttt{lbc.*.nc} reais de produção; e
funcionalmente, executando o \texttt{mpas\_atmosphere} real a partir
desses arquivos e obtendo uma previsão de 24 horas fisicamente sã, cuja
divergência frente à rota de referência cresce de forma suave e
limitada, consistente com o comportamento esperado de duas trajetórias
vizinhas de um sistema dinâmico caótico -- não com um artefato de
implementação. Essa validação, no entanto, foi conduzida em um único
par malha/caso de estudo (\cref{cap:discussao}); a rota foi generalizada
estruturalmente (leitura de \textit{namelist} real em tempo de execução)
para não depender de recompilação em outros cenários, mas ainda não foi
exercitada neles.

\section{Trabalhos futuros}

Três direções concretas de continuidade emergem diretamente das
limitações documentadas no \cref{cap:discussao}:

\begin{enumerate}
  \item \textbf{Implementar a mistura de terreno de fronteira}
    (\texttt{config\_blend\_bdy\_terrain}), fechando o resíduo residual
    de $\sim\SI{171}{\metre}$ concentrado no anel de células de
    fronteira -- a lacuna de maior impacto numérico mensurável
    remanescente, e a mais diretamente relevante caso a fronteira
    lateral se torne um ponto crítico de precisão em usos futuros.
  \item \textbf{Ampliar a validação para outras malhas, resoluções e
    estações do ano}, exercitando a generalização estrutural já
    implementada (\cref{sec:generalizacao}) e verificando se os 61
    níveis de pressão fixa intermediários -- derivados estatisticamente
    de uma única malha/caso -- permanecem adequados em contextos com
    perfil vertical de pressão substancialmente distinto.
  \item \textbf{Avaliar interpolantes de maior continuidade} sobre a
    malha dual de Delaunay (por exemplo, a família de interpolantes
    generalizados de Sibson/Laplace descrita por Hiyoshi e Sugihara
    \cite{hiyoshi2000continuity}) como substituto opcional da
    interpolação baricêntrica linear, caso alguma aplicação futura do
    mesmo motor de remapeamento -- por exemplo, fora do contexto de
    condição inicial/fronteira, em algum uso que dependa diretamente de
    gradiente horizontal suave -- se mostre sensível à descontinuidade
    de derivada nas arestas dos triângulos de Delaunay.
\end{enumerate}

Adicionalmente, uma investigação independente e de menor prioridade
identificada ao longo deste trabalho, mas fora de seu escopo direto: o
campo \texttt{lbc\_qv} dos arquivos de referência reais de produção
mostrou-se constante em todos os níveis verticais por célula
(\cref{cap:resultados}) -- um comportamento fisicamente implausível que
sugere um problema já existente na cadeia de geração da rota original
(via WPS), não introduzido por este trabalho, e que mereceria
investigação própria caso a fronteira lateral de referência continue
sendo usada como padrão de comparação em trabalhos futuros.


\clearpage
\printbibliography[heading=bibintoc,title={Referências}]

\end{document}
